\chapter{Concepción}\label{dept:Concepción}
\mapaparaguai{1}{Concepción}
\vfill\clearpage
\section{Concepción}

\subsection{Ficha climática de Concepción (capital
departamental)}

\textbf{Irradiância solar global --- (kWh/\allowbreak m²/\allowbreak dia):}

{\def\LTcaptype{none} % do not increment counter
\begin{longtable}[]{@{}
  >{\raggedright\arraybackslash}p{(\linewidth - 22\tabcolsep) * \real{0.0833}}
  >{\raggedright\arraybackslash}p{(\linewidth - 22\tabcolsep) * \real{0.0833}}
  >{\raggedright\arraybackslash}p{(\linewidth - 22\tabcolsep) * \real{0.0833}}
  >{\raggedright\arraybackslash}p{(\linewidth - 22\tabcolsep) * \real{0.0833}}
  >{\raggedright\arraybackslash}p{(\linewidth - 22\tabcolsep) * \real{0.0833}}
  >{\raggedright\arraybackslash}p{(\linewidth - 22\tabcolsep) * \real{0.0833}}
  >{\raggedright\arraybackslash}p{(\linewidth - 22\tabcolsep) * \real{0.0833}}
  >{\raggedright\arraybackslash}p{(\linewidth - 22\tabcolsep) * \real{0.0833}}
  >{\raggedright\arraybackslash}p{(\linewidth - 22\tabcolsep) * \real{0.0833}}
  >{\raggedright\arraybackslash}p{(\linewidth - 22\tabcolsep) * \real{0.0833}}
  >{\raggedright\arraybackslash}p{(\linewidth - 22\tabcolsep) * \real{0.0833}}
  >{\raggedright\arraybackslash}p{(\linewidth - 22\tabcolsep) * \real{0.0833}}@{}}
\toprule\noalign{}
\begin{minipage}[b]{\linewidth}\raggedright
Jan
\end{minipage} & \begin{minipage}[b]{\linewidth}\raggedright
Fev
\end{minipage} & \begin{minipage}[b]{\linewidth}\raggedright
Mar
\end{minipage} & \begin{minipage}[b]{\linewidth}\raggedright
Abr
\end{minipage} & \begin{minipage}[b]{\linewidth}\raggedright
Mai
\end{minipage} & \begin{minipage}[b]{\linewidth}\raggedright
Jun
\end{minipage} & \begin{minipage}[b]{\linewidth}\raggedright
Jul
\end{minipage} & \begin{minipage}[b]{\linewidth}\raggedright
Ago
\end{minipage} & \begin{minipage}[b]{\linewidth}\raggedright
Set
\end{minipage} & \begin{minipage}[b]{\linewidth}\raggedright
Out
\end{minipage} & \begin{minipage}[b]{\linewidth}\raggedright
Nov
\end{minipage} & \begin{minipage}[b]{\linewidth}\raggedright
Dez
\end{minipage} \\
\midrule\noalign{}
\endhead
\bottomrule\noalign{}
\endlastfoot
6.68 & 6.07 & 5.58 & 4.63 & 3.47 & 3.00 & 3.30 & 4.09 & 4.68 & 5.49 &
6.34 & 6.61 \\
\end{longtable}
}

Média anual: 4.99 kWh/\allowbreak m²/\allowbreak dia. Inclinação recomendada: 23° Norte (anual);
33° inverno; 13° verão.

\textbf{Precipitação --- (mm/\allowbreak mês):}

{\def\LTcaptype{none} % do not increment counter
\begin{longtable}[]{@{}
  >{\raggedright\arraybackslash}p{(\linewidth - 22\tabcolsep) * \real{0.0833}}
  >{\raggedright\arraybackslash}p{(\linewidth - 22\tabcolsep) * \real{0.0833}}
  >{\raggedright\arraybackslash}p{(\linewidth - 22\tabcolsep) * \real{0.0833}}
  >{\raggedright\arraybackslash}p{(\linewidth - 22\tabcolsep) * \real{0.0833}}
  >{\raggedright\arraybackslash}p{(\linewidth - 22\tabcolsep) * \real{0.0833}}
  >{\raggedright\arraybackslash}p{(\linewidth - 22\tabcolsep) * \real{0.0833}}
  >{\raggedright\arraybackslash}p{(\linewidth - 22\tabcolsep) * \real{0.0833}}
  >{\raggedright\arraybackslash}p{(\linewidth - 22\tabcolsep) * \real{0.0833}}
  >{\raggedright\arraybackslash}p{(\linewidth - 22\tabcolsep) * \real{0.0833}}
  >{\raggedright\arraybackslash}p{(\linewidth - 22\tabcolsep) * \real{0.0833}}
  >{\raggedright\arraybackslash}p{(\linewidth - 22\tabcolsep) * \real{0.0833}}
  >{\raggedright\arraybackslash}p{(\linewidth - 22\tabcolsep) * \real{0.0833}}@{}}
\toprule\noalign{}
\begin{minipage}[b]{\linewidth}\raggedright
Jan
\end{minipage} & \begin{minipage}[b]{\linewidth}\raggedright
Fev
\end{minipage} & \begin{minipage}[b]{\linewidth}\raggedright
Mar
\end{minipage} & \begin{minipage}[b]{\linewidth}\raggedright
Abr
\end{minipage} & \begin{minipage}[b]{\linewidth}\raggedright
Mai
\end{minipage} & \begin{minipage}[b]{\linewidth}\raggedright
Jun
\end{minipage} & \begin{minipage}[b]{\linewidth}\raggedright
Jul
\end{minipage} & \begin{minipage}[b]{\linewidth}\raggedright
Ago
\end{minipage} & \begin{minipage}[b]{\linewidth}\raggedright
Set
\end{minipage} & \begin{minipage}[b]{\linewidth}\raggedright
Out
\end{minipage} & \begin{minipage}[b]{\linewidth}\raggedright
Nov
\end{minipage} & \begin{minipage}[b]{\linewidth}\raggedright
Dez
\end{minipage} \\
\midrule\noalign{}
\endhead
\bottomrule\noalign{}
\endlastfoot
157 & 151 & 120 & 139 & 116 & 71 & 48 & 31 & 73 & 151 & 204 & 167 \\
\end{longtable}
}

Média anual: \textasciitilde1.427 mm/\allowbreak ano. Estação chuvosa Out-Mai; seca
Jun-Set.
\clearpage
\subsection*{Dados Departamentais}
\subsubsection{Desenvolvimento Humano e
Social}

\subsection{IDH Departamental}

{\def\LTcaptype{none} % do not increment counter
\begin{longtable}[]{@{}ll@{}}
\toprule\noalign{}
Indicador & Valor \\
\midrule\noalign{}
\endhead
\bottomrule\noalign{}
\endlastfoot
IDH & 0,700 \\
Classificação IDH & alto \\
Ranking nacional (1=melhor) & 9/\allowbreak 18 \\
Esperança de vida & 72,4 anos \\
Anos de escolaridade média & 8,1 anos \\
RNB per capita (USD PPA) & 8.400 \\
\end{longtable}
}

\subsection{Indicadores de Pobreza}

{\def\LTcaptype{none} % do not increment counter
\begin{longtable}[]{@{}ll@{}}
\toprule\noalign{}
Indicador & Valor \\
\midrule\noalign{}
\endhead
\bottomrule\noalign{}
\endlastfoot
Taxa de pobreza total (\%) & 31,1\% \\
Taxa de pobreza extrema (\%) & 7,8\% \\
Índice de Gini & 0,460 \\
\end{longtable}
}

\subsection{Infraestrutura Básica}

{\def\LTcaptype{none} % do not increment counter
\begin{longtable}[]{@{}ll@{}}
\toprule\noalign{}
Indicador & Valor \\
\midrule\noalign{}
\endhead
\bottomrule\noalign{}
\endlastfoot
Acesso a água potável (\%) & 74,2\% \\
Acesso a saneamento (\%) & 71,8\% \\
\end{longtable}
}

\subsubsection{Segurança Pública}

\subsection{Indicadores}

{\def\LTcaptype{none} % do not increment counter
\begin{longtable}[]{@{}
  >{\raggedright\arraybackslash}p{(\linewidth - 2\tabcolsep) * \real{0.6111}}
  >{\raggedright\arraybackslash}p{(\linewidth - 2\tabcolsep) * \real{0.3889}}@{}}
\toprule\noalign{}
\begin{minipage}[b]{\linewidth}\raggedright
Indicador
\end{minipage} & \begin{minipage}[b]{\linewidth}\raggedright
Valor
\end{minipage} \\
\midrule\noalign{}
\endhead
\bottomrule\noalign{}
\endlastfoot
Taxa de homicídios (por 100k hab) & \textasciitilde12--15 (estimativa,
acima da média nacional) \\
Crimes registrados (total/\allowbreak ano) & --- dados não desagregados
publicamente \\
Número de comisarías & \textasciitilde15 (estimativa departamental) \\
Índice segurança relativo & baixo \\
\end{longtable}
}

\subsubsection{Mercado de Terra Rural}

\subsection{Preços por Tipo de Terra}

{\def\LTcaptype{none} % do not increment counter
\begin{longtable}[]{@{}llll@{}}
\toprule\noalign{}
Tipo & Preço (USD/\allowbreak ha) & Faixa & Tendência \\
\midrule\noalign{}
\endhead
\bottomrule\noalign{}
\endlastfoot
Agrícola alta produtividade & 3.500 & 2.500 -- 5.000 & ↑ \\
Pastagem & 1.800 & 1.200 -- 2.500 & ↑ \\
Mata /\allowbreak  reserva & 600 & 300 -- 1.000 & → \\
\end{longtable}
}

\subsection{Características do
Mercado}

{\def\LTcaptype{none} % do not increment counter
\begin{longtable}[]{@{}
  >{\raggedright\arraybackslash}p{(\linewidth - 2\tabcolsep) * \real{0.6111}}
  >{\raggedright\arraybackslash}p{(\linewidth - 2\tabcolsep) * \real{0.3889}}@{}}
\toprule\noalign{}
\begin{minipage}[b]{\linewidth}\raggedright
Parâmetro
\end{minipage} & \begin{minipage}[b]{\linewidth}\raggedright
Valor
\end{minipage} \\
\midrule\noalign{}
\endhead
\bottomrule\noalign{}
\endlastfoot
Valorização anual média (\%) & 6--8\% \\
Lote mínimo rural (ha) & 20 \\
Imposto territorial rural (\%) & 1\% sobre valor fiscal (IRAGRO
adicional para renda) \\
Liquidez do mercado & média \\
\end{longtable}
}
\clearpage
\newpage
\secaoDiagnostico{Arroyito}{\mapaDistrito{0112}}{Arroyito é um distrito de origem camponesa com solo altamente produtivo,
mas marcado pela insegurança regional. Oferece baixo custo de vida e boa
capacidade de produção de alimentos, mas exige um alto nível de
tolerância ao risco devido à presença da insurgência no departamento.

\subsubsection{Por que sim}

\subsubsection{Por que não}

\subsubsection{Recomendação
Técnica}

Indicado para projetos de autossuficiência alimentar que possam ser
integrados à comunidade local e que possuam planos de segurança
robustos. Referência atual de score: GSS 5.5.

\begin{itemize}
\tightlist
\item
  \begin{enumerate}
  \def\labelenumi{\arabic{enumi})}
  \tightlist
  \item
    Dados censitários INE\footnote{\fineINE} 2022 (Arroyito Distrito).
  \end{enumerate}
\item
  \begin{enumerate}
  \def\labelenumi{\arabic{enumi})}
  \setcounter{enumi}{1}
  \tightlist
  \item
    Histórico de assentamento camponês (1989-2016).
  \end{enumerate}
\item
  \begin{enumerate}
  \def\labelenumi{\arabic{enumi})}
  \setcounter{enumi}{2}
  \tightlist
  \item
    Relatórios de segurança pública (FTC/\allowbreak CODI).
  \end{enumerate}
\item
  \begin{enumerate}
  \def\labelenumi{\arabic{enumi})}
  \setcounter{enumi}{3}
  \tightlist
  \item
    Lei de Segurança Fronteiriça (2532/\allowbreak 05).
  \end{enumerate}
\end{itemize}}{dist:01_Concepción:Arroyito}
\begin{table}[ht!]
\small \begin{tabular}{p{3.0cm}p{0.8cm}p{6.0cm}} \toprule
\textbf{Critério} & \textbf{Nota} & \textbf{Justificativa} \\ \midrule
A - Ameaças Estratégicas & 3.5 & Hotspot de insurgência regional e presença militar constante \\
B - Risco Social & 4.5 & Assentamento com histórico de agitação social e criminalidade rural \\
C - Riscos Naturais & 7.0 & Baixo risco de desastres geológicos; riscos climáticos manejáveis \\
D - Autossuficiência & 7.5 & Solo fértil e abundância de produção de alimentos básicos \\
E - Institucional & 5.5 & Ambiente institucional em consolidação; fortes restrições de fronteira \\
\midrule \textbf{GSS FINAL} & \textbf{5.5} & \textbf{Moderadamente Seguro (Recomendado apenas para perfis resilientes com foco em agricultura familiar e segurança tática).} \\ \bottomrule \end{tabular}
\end{table}
\subsubsection{Vulnerabilidades e
Mitigação}

\begin{itemize}
\tightlist
\item
  Exposicao hidrometeorologica moderada (45-64/\allowbreak 100).
\item
  Conectividade viaria intermediaria (50-69/\allowbreak 100).
\item
  Estoque de contingencia para 14 dias (agua/\allowbreak alimentos/\allowbreak combustivel),
  priorizado quando risco hidro=50/\allowbreak 100.
\item
  Duplicar rotas/\allowbreak logistica quando conectividade viaria=62/\allowbreak 100 for
  \textless70; manter rota secundaria mapeada e testada trimestralmente.
\item
  Plano de energia resiliente com autonomia minima de 72h
  (geracao/\allowbreak backup), revisao semestral.
\end{itemize}

\subsubsection{Dossiê de Campo}
\begin{description}[style=nextline,leftmargin=0.5cm,font=\bfseries]
\item[Coordenadas]  22°39'S 57°50'W (geocodificado via Nominatim).
\end{description}
\textbf{Fonte climática:} NASA POWER Climatology API (período 2001-2020)

\textbf{Fonte luminosa:} estimativa world atlas

\paragraph{Irradiação Solar
(kWh/\allowbreak m²/\allowbreak dia)}

{\def\LTcaptype{none} % do not increment counter
\begin{longtable}[]{@{}
  >{\raggedright\arraybackslash}p{(\linewidth - 24\tabcolsep) * \real{0.0746}}
  >{\raggedright\arraybackslash}p{(\linewidth - 24\tabcolsep) * \real{0.0746}}
  >{\raggedright\arraybackslash}p{(\linewidth - 24\tabcolsep) * \real{0.0746}}
  >{\raggedright\arraybackslash}p{(\linewidth - 24\tabcolsep) * \real{0.0746}}
  >{\raggedright\arraybackslash}p{(\linewidth - 24\tabcolsep) * \real{0.0746}}
  >{\raggedright\arraybackslash}p{(\linewidth - 24\tabcolsep) * \real{0.0746}}
  >{\raggedright\arraybackslash}p{(\linewidth - 24\tabcolsep) * \real{0.0746}}
  >{\raggedright\arraybackslash}p{(\linewidth - 24\tabcolsep) * \real{0.0746}}
  >{\raggedright\arraybackslash}p{(\linewidth - 24\tabcolsep) * \real{0.0746}}
  >{\raggedright\arraybackslash}p{(\linewidth - 24\tabcolsep) * \real{0.0746}}
  >{\raggedright\arraybackslash}p{(\linewidth - 24\tabcolsep) * \real{0.0746}}
  >{\raggedright\arraybackslash}p{(\linewidth - 24\tabcolsep) * \real{0.0746}}
  >{\raggedright\arraybackslash}p{(\linewidth - 24\tabcolsep) * \real{0.1045}}@{}}
\toprule\noalign{}
\begin{minipage}[b]{\linewidth}\raggedright
Jan
\end{minipage} & \begin{minipage}[b]{\linewidth}\raggedright
Fev
\end{minipage} & \begin{minipage}[b]{\linewidth}\raggedright
Mar
\end{minipage} & \begin{minipage}[b]{\linewidth}\raggedright
Abr
\end{minipage} & \begin{minipage}[b]{\linewidth}\raggedright
Mai
\end{minipage} & \begin{minipage}[b]{\linewidth}\raggedright
Jun
\end{minipage} & \begin{minipage}[b]{\linewidth}\raggedright
Jul
\end{minipage} & \begin{minipage}[b]{\linewidth}\raggedright
Ago
\end{minipage} & \begin{minipage}[b]{\linewidth}\raggedright
Set
\end{minipage} & \begin{minipage}[b]{\linewidth}\raggedright
Out
\end{minipage} & \begin{minipage}[b]{\linewidth}\raggedright
Nov
\end{minipage} & \begin{minipage}[b]{\linewidth}\raggedright
Dez
\end{minipage} & \begin{minipage}[b]{\linewidth}\raggedright
Média
\end{minipage} \\
\midrule\noalign{}
\endhead
\bottomrule\noalign{}
\endlastfoot
6.55 & 5.94 & 5.63 & 4.73 & 3.58 & 3.17 & 3.41 & 4.19 & 4.70 & 5.44 &
6.28 & 6.49 & \textbf{5.01} \\
\end{longtable}
}

\textbf{Inclinação solar recomendada:} 23° N (anual) · 33° N (inverno
jun-ago) · 13° N (verão nov-jan)

\paragraph{Precipitação (mm/\allowbreak mês)}

{\def\LTcaptype{none} % do not increment counter
\begin{longtable}[]{@{}
  >{\raggedright\arraybackslash}p{(\linewidth - 24\tabcolsep) * \real{0.0704}}
  >{\raggedright\arraybackslash}p{(\linewidth - 24\tabcolsep) * \real{0.0704}}
  >{\raggedright\arraybackslash}p{(\linewidth - 24\tabcolsep) * \real{0.0704}}
  >{\raggedright\arraybackslash}p{(\linewidth - 24\tabcolsep) * \real{0.0704}}
  >{\raggedright\arraybackslash}p{(\linewidth - 24\tabcolsep) * \real{0.0704}}
  >{\raggedright\arraybackslash}p{(\linewidth - 24\tabcolsep) * \real{0.0704}}
  >{\raggedright\arraybackslash}p{(\linewidth - 24\tabcolsep) * \real{0.0704}}
  >{\raggedright\arraybackslash}p{(\linewidth - 24\tabcolsep) * \real{0.0704}}
  >{\raggedright\arraybackslash}p{(\linewidth - 24\tabcolsep) * \real{0.0704}}
  >{\raggedright\arraybackslash}p{(\linewidth - 24\tabcolsep) * \real{0.0704}}
  >{\raggedright\arraybackslash}p{(\linewidth - 24\tabcolsep) * \real{0.0704}}
  >{\raggedright\arraybackslash}p{(\linewidth - 24\tabcolsep) * \real{0.0704}}
  >{\raggedright\arraybackslash}p{(\linewidth - 24\tabcolsep) * \real{0.1549}}@{}}
\toprule\noalign{}
\begin{minipage}[b]{\linewidth}\raggedright
Jan
\end{minipage} & \begin{minipage}[b]{\linewidth}\raggedright
Fev
\end{minipage} & \begin{minipage}[b]{\linewidth}\raggedright
Mar
\end{minipage} & \begin{minipage}[b]{\linewidth}\raggedright
Abr
\end{minipage} & \begin{minipage}[b]{\linewidth}\raggedright
Mai
\end{minipage} & \begin{minipage}[b]{\linewidth}\raggedright
Jun
\end{minipage} & \begin{minipage}[b]{\linewidth}\raggedright
Jul
\end{minipage} & \begin{minipage}[b]{\linewidth}\raggedright
Ago
\end{minipage} & \begin{minipage}[b]{\linewidth}\raggedright
Set
\end{minipage} & \begin{minipage}[b]{\linewidth}\raggedright
Out
\end{minipage} & \begin{minipage}[b]{\linewidth}\raggedright
Nov
\end{minipage} & \begin{minipage}[b]{\linewidth}\raggedright
Dez
\end{minipage} & \begin{minipage}[b]{\linewidth}\raggedright
Total/\allowbreak ano
\end{minipage} \\
\midrule\noalign{}
\endhead
\bottomrule\noalign{}
\endlastfoot
152 & 149 & 120 & 118 & 106 & 57 & 38 & 23 & 55 & 130 & 202 & 167 &
\textbf{1321 mm} \\
\end{longtable}
}

\paragraph{Poluição Luminosa}

{\def\LTcaptype{none} % do not increment counter
\begin{longtable}[]{@{}ll@{}}
\toprule\noalign{}
Parâmetro & Valor \\
\midrule\noalign{}
\endhead
\bottomrule\noalign{}
\endlastfoot
Escala Bortle & 2 --- Céu tipicamente escuro \\
Radiância artificial & 0.3 nW/\allowbreak cm²/\allowbreak sr \\
\end{longtable}
}
\paragraph{Solo (SoilGrids 2.0, média ponderada 0--30
cm)}

{\def\LTcaptype{none} % do not increment counter
\begin{longtable}[]{@{}ll@{}}
\toprule\noalign{}
Parâmetro & Valor \\
\midrule\noalign{}
\endhead
\bottomrule\noalign{}
\endlastfoot
pH (H$_2$O) & 6.15 \\
Carbono orgânico (SOC) & 18.97 g/\allowbreak kg \\
Argila & 27.17 \% \\
Areia & 34.14 \% \\
Silte (calc.) & 38.7 \% \\
Densidade aparente & 1.23 g/\allowbreak cm³ \\
\textbf{Aptidão agrícola} & \textbf{Alta} \\
\end{longtable}
}

\subsubsection{Indicadores Sociais}

\textbf{Fontes:} PNUD 2020, INE\footnote{\fineINE} Censo 2022, INE\footnote{\fineINE} EPHC 2023, Ministerio
Público 2024\\
\textbf{Nota:} Valores marcados como \emph{(dept.)} referem-se ao
departamento de Presidente Hayes; valores \emph{(dist.)} são específicos
deste distrito. Dados marcados \emph{(est.)} são estimativas.

{\def\LTcaptype{none} % do not increment counter
\begin{longtable}[]{@{}
  >{\raggedright\arraybackslash}p{(\linewidth - 4\tabcolsep) * \real{0.4231}}
  >{\raggedright\arraybackslash}p{(\linewidth - 4\tabcolsep) * \real{0.2692}}
  >{\raggedright\arraybackslash}p{(\linewidth - 4\tabcolsep) * \real{0.3077}}@{}}
\toprule\noalign{}
\begin{minipage}[b]{\linewidth}\raggedright
Indicador
\end{minipage} & \begin{minipage}[b]{\linewidth}\raggedright
Valor
\end{minipage} & \begin{minipage}[b]{\linewidth}\raggedright
Âmbito
\end{minipage} \\
\midrule\noalign{}
\endhead
\bottomrule\noalign{}
\endlastfoot
IDH (2020) & 0,693 & dept. \\
Ranking IDH nacional & 13/\allowbreak 18 & dept. \\
Esperança de vida & 71,8 anos & dept. \\
Escolaridade média & 4.7 anos & dist. \\
RNB per capita (USD PPA) & 7.900 & dept. \\
Pobreza monetária (\%) & 29,4\% & dept. \\
Pobreza extrema (\%) & 7,2\% & dept. \\
Índice de Gini & 0,453 & dept. \\
Acesso a água potável (\%) & 71,3\% & dept. \\
Acesso a saneamento (\%) & 68,9\% & dept. \\
Taxa de homicídios (est., /\allowbreak 100k hab) & 5,5--5,6 (2019--2020, fonte INE
--- próxima à média nacional) & dept. \\
Índice de segurança & médio & dept. \\
População (Censo 2022) & N/\allowbreak D & dist. \\
Idade mediana & N/\allowbreak D & dist. \\
Taxa de fecundidade & N/\allowbreak D & dist. \\
\end{longtable}
}
\subsubsection{Combustível}

\textbf{Referência:} PETROPAR /\allowbreak  postos locais (2024) \textbf{Tipo de
localidade:} Interior

{\def\LTcaptype{none} % do not increment counter
\begin{longtable}[]{@{}lll@{}}
\toprule\noalign{}
Combustível & USD/\allowbreak litro & Gs/\allowbreak litro (aprox.) \\
\midrule\noalign{}
\endhead
\bottomrule\noalign{}
\endlastfoot
Gasolina 93 oct & 1.05 & 7,770 \\
Gasolina 97 oct (premium) & 1.15 & 8,510 \\
Diesel & 0.98 & 7,252 \\
\end{longtable}
}

\begin{quote}
Preços podem variar ±5\% conforme posto e sazonalidade. Chaco e interior
remoto apresentam maior variação.
\end{quote}

\subsubsection{Cobertura Celular}

\textbf{Fonte:} CONATEL PY /\allowbreak  operadoras (2024)

{\def\LTcaptype{none} % do not increment counter
\begin{longtable}[]{@{}ll@{}}
\toprule\noalign{}
Parâmetro & Valor \\
\midrule\noalign{}
\endhead
\bottomrule\noalign{}
\endlastfoot
Cobertura 4G população (dept.) & 60\% \\
Cobertura 4G área rural & 35\% \\
Melhor operadora & Tigo \\
Qualidade rural & limitada \\
\end{longtable}
}

\begin{quote}
Para áreas rurais fora do núcleo urbano, recomenda-se chip Tigo como
principal e Personal como backup.
\end{quote}

\subsubsection{Internet}

\textbf{Fonte:} CONATEL /\allowbreak  Speedtest Ookla (2024)

{\def\LTcaptype{none} % do not increment counter
\begin{longtable}[]{@{}ll@{}}
\toprule\noalign{}
Parâmetro & Valor \\
\midrule\noalign{}
\endhead
\bottomrule\noalign{}
\endlastfoot
Velocidade média download & 25 Mbps \\
Domicílios com internet (dept.) & 35\% \\
Tecnologia predominante & rádio/\allowbreak satélite \\
Opção rural & Starlink disponível (\textasciitilde USD 44/\allowbreak mês) \\
\end{longtable}
}

\subsubsection{Mercado Imobiliário e Terra
Rural}

\textbf{Fonte:} INDERT /\allowbreak  Clasificados.com.py (2024)

{\def\LTcaptype{none} % do not increment counter
\begin{longtable}[]{@{}ll@{}}
\toprule\noalign{}
Tipo & Referência \\
\midrule\noalign{}
\endhead
\bottomrule\noalign{}
\endlastfoot
Terra agrícola alta prod. (USD/\allowbreak ha) & 2,000 \\
Imóvel urbano (USD/\allowbreak m²) & 600 \\
Aluguel 2 quartos (USD/\allowbreak mês) & 250 \\
\end{longtable}
}

\begin{quote}
Valores de referência departamental. Localidades menores podem ter
preços 20--40\% abaixo da capital departamental.
\end{quote}

\subsubsection{Saúde}

\textbf{Fonte:} MSPBS /\allowbreak  IPS Paraguay (2024-2026), consolidação
departamental e proxy local

{\def\LTcaptype{none} % do not increment counter
\begin{longtable}[]{@{}ll@{}}
\toprule\noalign{}
Serviço & Disponibilidade \\
\midrule\noalign{}
\endhead
\bottomrule\noalign{}
\endlastfoot
USF /\allowbreak  Posto de Saúde & sim \\
Hospital Regional & não \\
IPS (seguro social) & não \\
Farmácia & sim \\
Distância ao hospital de referência & \textasciitilde343 km (Villa
Hayes) \\
\end{longtable}
}

\textbf{Principais estabelecimentos:} USF local; referencia hospitalar
em Villa Hayes

\textbf{Observação para imigrantes:} Atendimento primario local ou em
raio curto; casos de maior complexidade seguem para Villa Hayes.
Cobertura privada continua recomendada para especialidades e urgências
de maior complexidade.
\newpage
\secaoDiagnostico{Azotey}{\mapaConteudo{-23.2166}{-56.5000}}{Azotey é um ponto estratégico de controle rodoviário no Norte,
oferecendo terras de alta aptidão pecuária e agrícola. Contudo, seu GSS
de 5.3 reflete a realidade de uma zona de conflito ativo onde a
segurança tática deve ser a prioridade número um de qualquer ocupante.

\subsubsection{Por que sim}

\subsubsection{Por que não}

\subsubsection{Recomendação
Técnica}

Não recomendado para residência civil passiva. Indicado apenas como
ativo de investimento agroindustrial para grupos que possuam estrutura
de segurança e contingência profissional. Referência atual de score: GSS
5.3.

\begin{itemize}
\tightlist
\item
  \begin{enumerate}
  \def\labelenumi{\arabic{enumi})}
  \tightlist
  \item
    Estimativa populacional distrital (INE\footnote{\fineINE} 2022).
  \end{enumerate}
\item
  \begin{enumerate}
  \def\labelenumi{\arabic{enumi})}
  \setcounter{enumi}{1}
  \tightlist
  \item
    Relatórios de conflitos camponeses e guerrilha (FTC 2020-2025).
  \end{enumerate}
\item
  \begin{enumerate}
  \def\labelenumi{\arabic{enumi})}
  \setcounter{enumi}{2}
  \tightlist
  \item
    Dados de mercado de terras no Norte paraguaio.
  \end{enumerate}
\item
  \begin{enumerate}
  \def\labelenumi{\arabic{enumi})}
  \setcounter{enumi}{3}
  \tightlist
  \item
    Delimitação de faixa de segurança de fronteira.
  \end{enumerate}
\end{itemize}}{dist:01_Concepción:Azotey}
\begin{table}[ht!]
\small \begin{tabular}{p{3.0cm}p{0.8cm}p{6.0cm}} \toprule
\textbf{Critério} & \textbf{Nota} & \textbf{Justificativa} \\ \midrule
A - Ameaças Estratégicas & 3.0 & Alvo tático de grupos armados irregulares; presença militar intensa \\
B - Risco Social & 4.0 & Risco social elevado devido a conflitos de terra e guerrilha rural \\
C - Riscos Naturais & 7.0 & Geologia estável; clima favorável à produção mas com isolamento pluvial \\
D - Autossuficiência & 7.5 & Solo produtivo e abundância de proteína animal em estâncias vizinhas \\
E - Institucional & 5.5 & Ambiente institucional sob tensão de segurança; restrições legais para estrangeiros \\
\midrule \textbf{GSS FINAL} & \textbf{5.3} & \textbf{Sensivel (Recomendado apenas para perfis táticos de alta resiliência e interesse específico em agronegócio protegido).} \\ \bottomrule \end{tabular}
\end{table}
\paragraph{Solo (SoilGrids 2.0, média ponderada 0--30
cm)}

{\def\LTcaptype{none} % do not increment counter
\begin{longtable}[]{@{}ll@{}}
\toprule\noalign{}
Parâmetro & Valor \\
\midrule\noalign{}
\endhead
\bottomrule\noalign{}
\endlastfoot
pH (H$_2$O) & 6.18 \\
Carbono orgânico (SOC) & 17.32 g/\allowbreak kg \\
Argila & 38.34 \% \\
Areia & 20.41 \% \\
Silte (calc.) & 41.2 \% \\
Densidade aparente & 1.33 g/\allowbreak cm³ \\
\textbf{Aptidão agrícola} & \textbf{Alta} \\
\end{longtable}
}

\subsubsection{Indicadores Sociais}

\textbf{Fontes:} PNUD 2020, INE\footnote{\fineINE} Censo 2022, INE\footnote{\fineINE} EPHC 2023, Ministerio
Público 2024\\
\textbf{Nota:} Valores marcados como \emph{(dept.)} referem-se ao
departamento de Presidente Hayes; valores \emph{(dist.)} são específicos
deste distrito. Dados marcados \emph{(est.)} são estimativas.

{\def\LTcaptype{none} % do not increment counter
\begin{longtable}[]{@{}
  >{\raggedright\arraybackslash}p{(\linewidth - 4\tabcolsep) * \real{0.4231}}
  >{\raggedright\arraybackslash}p{(\linewidth - 4\tabcolsep) * \real{0.2692}}
  >{\raggedright\arraybackslash}p{(\linewidth - 4\tabcolsep) * \real{0.3077}}@{}}
\toprule\noalign{}
\begin{minipage}[b]{\linewidth}\raggedright
Indicador
\end{minipage} & \begin{minipage}[b]{\linewidth}\raggedright
Valor
\end{minipage} & \begin{minipage}[b]{\linewidth}\raggedright
Âmbito
\end{minipage} \\
\midrule\noalign{}
\endhead
\bottomrule\noalign{}
\endlastfoot
IDH (2020) & 0,693 & dept. \\
Ranking IDH nacional & 13/\allowbreak 18 & dept. \\
Esperança de vida & 71,8 anos & dept. \\
Escolaridade média & 7.4 anos & dist. \\
RNB per capita (USD PPA) & 7.900 & dept. \\
Pobreza monetária (\%) & 29,4\% & dept. \\
Pobreza extrema (\%) & 7,2\% & dept. \\
Índice de Gini & 0,453 & dept. \\
Acesso a água potável (\%) & 71,3\% & dept. \\
Acesso a saneamento (\%) & 68,9\% & dept. \\
Taxa de homicídios (est., /\allowbreak 100k hab) & 5,5--5,6 (2019--2020, fonte INE
--- próxima à média nacional) & dept. \\
Índice de segurança & médio & dept. \\
População (Censo 2022) & 2.743 habitantes & dist. \\
Idade mediana & N/\allowbreak D & dist. \\
Taxa de fecundidade & N/\allowbreak D & dist. \\
\end{longtable}
}
\subsubsection{Combustível}

\textbf{Referência:} PETROPAR /\allowbreak  postos locais (2024) \textbf{Tipo de
localidade:} Interior

{\def\LTcaptype{none} % do not increment counter
\begin{longtable}[]{@{}lll@{}}
\toprule\noalign{}
Combustível & USD/\allowbreak litro & Gs/\allowbreak litro (aprox.) \\
\midrule\noalign{}
\endhead
\bottomrule\noalign{}
\endlastfoot
Gasolina 93 oct & 1.05 & 7,770 \\
Gasolina 97 oct (premium) & 1.15 & 8,510 \\
Diesel & 0.98 & 7,252 \\
\end{longtable}
}

\begin{quote}
Preços podem variar ±5\% conforme posto e sazonalidade. Chaco e interior
remoto apresentam maior variação.
\end{quote}

\subsubsection{Cobertura Celular}

\textbf{Fonte:} CONATEL PY /\allowbreak  operadoras (2024)

{\def\LTcaptype{none} % do not increment counter
\begin{longtable}[]{@{}ll@{}}
\toprule\noalign{}
Parâmetro & Valor \\
\midrule\noalign{}
\endhead
\bottomrule\noalign{}
\endlastfoot
Cobertura 4G população (dept.) & 60\% \\
Cobertura 4G área rural & 35\% \\
Melhor operadora & Tigo \\
Qualidade rural & limitada \\
\end{longtable}
}

\begin{quote}
Para áreas rurais fora do núcleo urbano, recomenda-se chip Tigo como
principal e Personal como backup.
\end{quote}

\subsubsection{Internet}

\textbf{Fonte:} CONATEL /\allowbreak  Speedtest Ookla (2024)

{\def\LTcaptype{none} % do not increment counter
\begin{longtable}[]{@{}ll@{}}
\toprule\noalign{}
Parâmetro & Valor \\
\midrule\noalign{}
\endhead
\bottomrule\noalign{}
\endlastfoot
Velocidade média download & 25 Mbps \\
Domicílios com internet (dept.) & 35\% \\
Tecnologia predominante & rádio/\allowbreak satélite \\
Opção rural & Starlink disponível (\textasciitilde USD 44/\allowbreak mês) \\
\end{longtable}
}

\subsubsection{Mercado Imobiliário e Terra
Rural}

\textbf{Fonte:} INDERT /\allowbreak  Clasificados.com.py (2024)

{\def\LTcaptype{none} % do not increment counter
\begin{longtable}[]{@{}ll@{}}
\toprule\noalign{}
Tipo & Referência \\
\midrule\noalign{}
\endhead
\bottomrule\noalign{}
\endlastfoot
Terra agrícola alta prod. (USD/\allowbreak ha) & 2,000 \\
Imóvel urbano (USD/\allowbreak m²) & 600 \\
Aluguel 2 quartos (USD/\allowbreak mês) & 250 \\
\end{longtable}
}

\begin{quote}
Valores de referência departamental. Localidades menores podem ter
preços 20--40\% abaixo da capital departamental.
\end{quote}

\subsubsection{Saúde}

\textbf{Fonte:} MSPBS /\allowbreak  IPS Paraguay (2024-2026), consolidação
departamental e proxy local

{\def\LTcaptype{none} % do not increment counter
\begin{longtable}[]{@{}ll@{}}
\toprule\noalign{}
Serviço & Disponibilidade \\
\midrule\noalign{}
\endhead
\bottomrule\noalign{}
\endlastfoot
USF /\allowbreak  Posto de Saúde & sim \\
Hospital Regional & não \\
IPS (seguro social) & não \\
Farmácia & sim \\
Distância ao hospital de referência & \textasciitilde322 km (Villa
Hayes) \\
\end{longtable}
}

\textbf{Principais estabelecimentos:} USF local; referencia hospitalar
em Villa Hayes

\textbf{Observação para imigrantes:} Atendimento primario local ou em
raio curto; casos de maior complexidade seguem para Villa Hayes.
Cobertura privada continua recomendada para especialidades e urgências
de maior complexidade.
\newpage
\secaoDiagnostico{Belen}{\mapaDistrito{0102}}{Belén é um polo de segurança alimentar e histórica no Norte. Sua posição
no Trópico de Capricórnio e a presença de uma das indústrias cárneas
mais modernas da região fazem dela um ativo estratégico para suprimento
e autossuficiência, com o bônus da liberdade jurídica para proprietários
estrangeiros.

\subsubsection{Por que sim}

\begin{itemize}
\tightlist
\item
  Hub de proteína animal de elite e abundância hídrica.
\item
  Fora da restrição legal de 50km da fronteira.
\item
  Proximidade imediata com a infraestrutura de Concepción.
\end{itemize}

\subsubsection{Por que não}

\subsubsection{Recomendação
Técnica}

Recomendado para perfis que buscam integrar investimento produtivo
(pecuária/\allowbreak reflorestamento) com segurança alimentar robusta. Referência
atual de score: GSS 6.4.

\begin{itemize}
\tightlist
\item
  \begin{enumerate}
  \def\labelenumi{\arabic{enumi})}
  \tightlist
  \item
    Dados censitários INE\footnote{\fineINE} 2022.
  \end{enumerate}
\item
  \begin{enumerate}
  \def\labelenumi{\arabic{enumi})}
  \setcounter{enumi}{1}
  \tightlist
  \item
    Relatórios de indústria frigorífica regional.
  \end{enumerate}
\item
  \begin{enumerate}
  \def\labelenumi{\arabic{enumi})}
  \setcounter{enumi}{2}
  \tightlist
  \item
    Mapa climático e geográfico (Trópico de Capricórnio).
  \end{enumerate}
\item
  \begin{enumerate}
  \def\labelenumi{\arabic{enumi})}
  \setcounter{enumi}{3}
  \tightlist
  \item
    Certidão de localização fora da faixa de fronteira.
  \end{enumerate}
\end{itemize}}{dist:01_Concepción:Belen}
\begin{table}[ht!]
\small \begin{tabular}{p{3.0cm}p{0.8cm}p{6.0cm}} \toprule
\textbf{Critério} & \textbf{Nota} & \textbf{Justificativa} \\ \midrule
A - Ameaças Estratégicas & 5.5 & Centro geográfico regional; alvo logístico de suprimentos proteicos \\
B - Risco Social & 6.0 & Densidade urbana manejável; controle policial tático permanente \\
C - Riscos Naturais & 6.5 & Geologia estável; risco hídrico moderado ligado ao Rio Ypané \\
D - Autossuficiência & 8.0 & Abundância extrema de proteína animal, água e solo fértil para erva-mate \\
E - Institucional & 7.0 & Ambiente comercial dinâmico; liberdade de aquisição para estrangeiros \\
\midrule \textbf{GSS FINAL} & \textbf{6.4} & \textbf{Moderadamente Seguro (Excelente opção para bases produtivas e autossuficiência agroindustrial).} \\ \bottomrule \end{tabular}
\end{table}
\subsubsection{Vulnerabilidades e
Mitigação}

\begin{itemize}
\tightlist
\item
  \textbf{Dependência Industrial:} Economia local muito vinculada ao
  frigorífico (Mitigação: diversificação para agricultura de
  subsistência e energia solar).
\item
  \textbf{Instabilidade Climática:} Temperaturas de verão exigem
  construção com alta inércia térmica e sistemas de resfriamento
  eficientes.
\item
  \textbf{Segurança Regional:} Monitoramento de movimentos insurgentes
  em distritos vizinhos (Mitigação: residência próxima ao perímetro
  urbano e comunicação satelital).
\end{itemize}

\subsubsection{Dossiê de Campo}
\begin{description}[style=nextline,leftmargin=0.5cm,font=\bfseries]
\item[Coordenadas]  23°27'S 57°15'W (Ponto exato de passagem do Trópico de Capricórnio).
\item[Topografia]  Relevo predominantemente plano com áreas de várzea próximas ao Rio Ypané. Altitude \textasciitilde 80m. Geologicamente estável, risco sísmico desprezível.
\item[Alvos Estratégicos]  Frigorífico Athena Foods (ex-JBS) - motor econômico e polo de suprimento proteico; Cruzamento de rotas para Concepción e San Pedro.
\item[Fallout]  Ventos predominantes do Norte, Leste e Sudeste.
\item[População]  10.605 habitantes (Censo 2022 Final).
\item[Densidade]  \textasciitilde 37,2 hab/\allowbreak km² (Moderada, com núcleo urbano consolidado e histórico).
\item[Vias de Acesso]  Conectividade asfáltica eficiente com Concepción (21 km). Ponto de trânsito regional de carga pesada (gado e grãos).
\item[Serviços]  Centro de saúde local para urgências básicas. Dependência direta do Hospital Regional de Concepción para alta complexidade.
\item[Custo de Vida]  Significativamente baixo (30-40\% menor que o Brasil para alimentação).
\item[Preço da Terra]  US\$ 4.000-10.000/\allowbreak ha (terras de alta produtividade agrícola/\allowbreak pastagem); US\$ 2.500-3.500/\allowbreak ha (áreas de reflorestamento).
\item[Hidrologia]  Risco Moderado de Inundação. Presença do Rio Ypané; áreas ribeirinhas vulneráveis em anos de El Niño intenso. Praias fluviais são ativos de lazer.
\item[Clima]  Tropical. Temperaturas extremas (máximas de 45°C e geadas ocasionais no inverno). Chuvas concentradas entre novembro e janeiro.
\item[Energia]  Rede nacional (Itaipu). Estabilidade moderada para o núcleo industrial.
\item[Água]  Abundante (Rio Ypané e poços artesianos).
\item[Qualidade do Solo]  Solos podzólicos e lateríticos; profundos e ácidos, respondem muito bem a fertilizantes. Ideais para pastagens, erva-mate e fruticultura.
\item[Recursos Locais]  Grande excedente de proteína animal (frigorífico local). Elevadíssimo potencial para autossuficiência e suprimento de longo prazo.
\item[Segurança]  Zona de atuação da FTC (Força de Tarefa Conjunta) devido à proximidade com áreas de insurgência (EPP\footnote{\fineEPP}) no departamento. Segurança urbana estável, policiamento ostensivo em rotas industriais. Taxa de homicídios regional de \textasciitilde 14/\allowbreak 100k.
\item[Leis Local: Município histórico (fundado em 1760). Livre de Restrição de Fronteira]  Localizado fora da faixa de 50km da Lei 2532/\allowbreak 05, facilitando a aquisição de terras rurais por estrangeiros brasileiros.
\end{description}
\subsubsection{Indicadores Sociais}

\textbf{Fontes:} PNUD 2020, INE\footnote{\fineINE} Censo 2022, INE\footnote{\fineINE} EPHC 2023, Ministerio
Público 2024\\
\textbf{Nota:} Valores marcados como \emph{(dept.)} referem-se ao
departamento de Presidente Hayes; valores \emph{(dist.)} são específicos
deste distrito. Dados marcados \emph{(est.)} são estimativas.

{\def\LTcaptype{none} % do not increment counter
\begin{longtable}[]{@{}
  >{\raggedright\arraybackslash}p{(\linewidth - 4\tabcolsep) * \real{0.4231}}
  >{\raggedright\arraybackslash}p{(\linewidth - 4\tabcolsep) * \real{0.2692}}
  >{\raggedright\arraybackslash}p{(\linewidth - 4\tabcolsep) * \real{0.3077}}@{}}
\toprule\noalign{}
\begin{minipage}[b]{\linewidth}\raggedright
Indicador
\end{minipage} & \begin{minipage}[b]{\linewidth}\raggedright
Valor
\end{minipage} & \begin{minipage}[b]{\linewidth}\raggedright
Âmbito
\end{minipage} \\
\midrule\noalign{}
\endhead
\bottomrule\noalign{}
\endlastfoot
IDH (2020) & 0,693 & dept. \\
Ranking IDH nacional & 13/\allowbreak 18 & dept. \\
Esperança de vida & 71,8 anos & dept. \\
Escolaridade média & 7,8 anos & dept. \\
RNB per capita (USD PPA) & 7.900 & dept. \\
Pobreza monetária (\%) & 29,4\% & dept. \\
Pobreza extrema (\%) & 7,2\% & dept. \\
Índice de Gini & 0,453 & dept. \\
Acesso a água potável (\%) & 71,3\% & dept. \\
Acesso a saneamento (\%) & 68,9\% & dept. \\
Taxa de homicídios (est., /\allowbreak 100k hab) & 5,5--5,6 (2019--2020, fonte INE
--- próxima à média nacional) & dept. \\
Índice de segurança & médio & dept. \\
População (Censo 2022) & 20.145 habitantes & dist. \\
Idade mediana & N/\allowbreak D & dist. \\
Taxa de fecundidade & N/\allowbreak D & dist. \\
\end{longtable}
}
\subsubsection{Combustível}

\textbf{Referência:} PETROPAR /\allowbreak  postos locais (2024) \textbf{Tipo de
localidade:} Interior

{\def\LTcaptype{none} % do not increment counter
\begin{longtable}[]{@{}lll@{}}
\toprule\noalign{}
Combustível & USD/\allowbreak litro & Gs/\allowbreak litro (aprox.) \\
\midrule\noalign{}
\endhead
\bottomrule\noalign{}
\endlastfoot
Gasolina 93 oct & 1.05 & 7,770 \\
Gasolina 97 oct (premium) & 1.15 & 8,510 \\
Diesel & 0.98 & 7,252 \\
\end{longtable}
}

\begin{quote}
Preços podem variar ±5\% conforme posto e sazonalidade. Chaco e interior
remoto apresentam maior variação.
\end{quote}

\subsubsection{Cobertura Celular}

\textbf{Fonte:} CONATEL PY /\allowbreak  operadoras (2024)

{\def\LTcaptype{none} % do not increment counter
\begin{longtable}[]{@{}ll@{}}
\toprule\noalign{}
Parâmetro & Valor \\
\midrule\noalign{}
\endhead
\bottomrule\noalign{}
\endlastfoot
Cobertura 4G população (dept.) & 60\% \\
Cobertura 4G área rural & 35\% \\
Melhor operadora & Tigo \\
Qualidade rural & limitada \\
\end{longtable}
}

\begin{quote}
Para áreas rurais fora do núcleo urbano, recomenda-se chip Tigo como
principal e Personal como backup.
\end{quote}

\subsubsection{Internet}

\textbf{Fonte:} CONATEL /\allowbreak  Speedtest Ookla (2024)

{\def\LTcaptype{none} % do not increment counter
\begin{longtable}[]{@{}ll@{}}
\toprule\noalign{}
Parâmetro & Valor \\
\midrule\noalign{}
\endhead
\bottomrule\noalign{}
\endlastfoot
Velocidade média download & 25 Mbps \\
Domicílios com internet (dept.) & 35\% \\
Tecnologia predominante & rádio/\allowbreak satélite \\
Opção rural & Starlink disponível (\textasciitilde USD 44/\allowbreak mês) \\
\end{longtable}
}

\subsubsection{Mercado Imobiliário e Terra
Rural}

\textbf{Fonte:} INDERT /\allowbreak  Clasificados.com.py (2024)

{\def\LTcaptype{none} % do not increment counter
\begin{longtable}[]{@{}ll@{}}
\toprule\noalign{}
Tipo & Referência \\
\midrule\noalign{}
\endhead
\bottomrule\noalign{}
\endlastfoot
Terra agrícola alta prod. (USD/\allowbreak ha) & 2,000 \\
Imóvel urbano (USD/\allowbreak m²) & 600 \\
Aluguel 2 quartos (USD/\allowbreak mês) & 250 \\
\end{longtable}
}

\begin{quote}
Valores de referência departamental. Localidades menores podem ter
preços 20--40\% abaixo da capital departamental.
\end{quote}

\subsubsection{Saúde}

\textbf{Fonte:} MSPBS /\allowbreak  IPS Paraguay (2024-2026), consolidação
departamental e proxy local

{\def\LTcaptype{none} % do not increment counter
\begin{longtable}[]{@{}ll@{}}
\toprule\noalign{}
Serviço & Disponibilidade \\
\midrule\noalign{}
\endhead
\bottomrule\noalign{}
\endlastfoot
USF /\allowbreak  Posto de Saúde & sim \\
Hospital Regional & sim \\
IPS (seguro social) & não \\
Farmácia & sim \\
Distância ao hospital de referência & \textasciitilde387 km (Villa
Hayes) \\
\end{longtable}
}

\textbf{Principais estabelecimentos:} USF local; referencia hospitalar
em Villa Hayes

\textbf{Observação para imigrantes:} Atendimento primario local ou em
raio curto; casos de maior complexidade seguem para Villa Hayes.
Cobertura privada continua recomendada para especialidades e urgências
de maior complexidade.
\newpage
\secaoDiagnostico{Concepción}{\mapaDistrito{0101}}{Concepción é o nó logístico e militar do Norte paraguaio. Oferece
recursos naturais excepcionais (solo e água), mas a pontuação GSS é
penalizada pelos riscos de segurança regional e vulnerabilidade a
cheias.

\subsubsection{Por que sim}

\subsubsection{Por que não}

\begin{itemize}
\tightlist
\item
  Insegurança regional (EPP\footnote{\fineEPP}/\allowbreak Narcotráfico) e criminalidade urbana mais
  alta.
\item
  Risco hidrológico severo.
\end{itemize}

\subsubsection{Recomendação
Técnica}

Indicado para perfis que buscam exploração agroindustrial de alto
desempenho, cientes da necessidade de investimento robusto em segurança
e mitigação hídrica. Referência atual de score: GSS 5.7.

\begin{itemize}
\tightlist
\item
  \begin{enumerate}
  \def\labelenumi{\arabic{enumi})}
  \tightlist
  \item
    Dados censitários (INE\footnote{\fineINE} 2022).
  \end{enumerate}
\item
  \begin{enumerate}
  \def\labelenumi{\arabic{enumi})}
  \setcounter{enumi}{1}
  \tightlist
  \item
    Relatórios de segurança regional (CODI 2025).
  \end{enumerate}
\item
  \begin{enumerate}
  \def\labelenumi{\arabic{enumi})}
  \setcounter{enumi}{2}
  \tightlist
  \item
    Mapas de uso do solo e sismicidade (Embrapa/\allowbreak ResearchGate).
  \end{enumerate}
\end{itemize}}{dist:01_Concepción:Concepción}
\begin{table}[ht!]
\small \begin{tabular}{p{3.0cm}p{0.8cm}p{6.0cm}} \toprule
\textbf{Critério} & \textbf{Nota} & \textbf{Justificativa} \\ \midrule
A - Ameaças Estratégicas & 4.0 & Base militar estratégica e centro de operações de defesa interna \\
B - Risco Social & 4.5 & Instabilidade regional por insurgência e criminalidade acima da média \\
C - Riscos Naturais & 6.0 & Geologia estável; risco crítico de inundação fluvial \\
D - Autossuficiência & 8.0 & Solo de alta produtividade e abundância hídrica/\allowbreak proteica \\
E - Institucional & 6.5 & Forte presença militar; restrições de fronteira para estrangeiros \\
\midrule \textbf{GSS FINAL} & \textbf{5.7} & \textbf{Moderadamente Seguro (Exige atenção rigorosa à segurança e riscos hídricos).} \\ \bottomrule \end{tabular}
\end{table}
\subsubsection{Vulnerabilidades e
Mitigação}

\begin{itemize}
\tightlist
\item
  \textbf{Instabilidade Regional:} Presença de grupos armados no
  departamento (Mitigação: residência no perímetro urbano controlado ou
  segurança privada dedicada em zonas rurais).
\item
  \textbf{Risco Hídrico:} Inundações severas (Mitigação: evitar áreas de
  cota baixa, mapear rotas de saída para o Chaco ou Sul).
\item
  \textbf{Pico Logístico:} Dependência de poucas pontes para cruzar o
  Rio Paraguai.
\end{itemize}

\subsubsection{Dossiê de Campo}
\begin{description}[style=nextline,leftmargin=0.5cm,font=\bfseries]
\item[Coordenadas]  23°24'S 57°26'W.
\item[Topografia]  Planície próxima ao Rio Paraguai. Altitude \textasciitilde 70m. Região tectonicamente estável, risco sísmico quase nulo.
\item[Alvos Estratégicos]  Porto de Concepción (logística de grãos e cimento), Aeroporto ANC (em modernização), Sede do CODI (Comando de Operações de Defesa Interna), Regimento de Infantaria N° 10 "Sauce". Planta da Paracel (Celulose) em implantação.
\item[Fallout]  Ventos predominantes do Norte/\allowbreak Nordeste (quentes) e Sul (frios no inverno).
\item[População]  73.360 (Censo 2022 Final).
\item[Densidade]  \textasciitilde 85 hab/\allowbreak km² (distrito total), zona urbana densa.
\item[Vias de Acesso]  Ruta PY05 (Conexão Leste-Oeste), Ponte sobre o Rio Paraguai. Gargalos logísticos devido à dependência de poucas travessias fluviais.
\item[Serviços]  Gran Hospital de Concepción (em construção, referência para o Norte). Sede administrativa regional.
\item[Custo de Vida]  US\$ 800-1.200/\allowbreak mês. Lotes urbanos: US\$ 15-30/\allowbreak m². Terras rurais: US\$ 5.000-8.000/\allowbreak ha.
\item[Hidrologia]  Risco Crítico de Inundação. Bairros ribeirinhos vulneráveis a cheias do Rio Paraguai e transbordamentos súbitos do Rio Ypané.
\item[Clima]  Tropical Úmido. Verões com temperaturas extremas. Tempestades intensas frequentes entre outubro e maio.
\item[Inclinação recomendada para placas solares]  23° voltado para o Norte (anual); 33° no inverno (Jun-Ago); 13° no verão (Nov-Jan). Base: latitude local -23.40°S.
\item[Energia]  Rede nacional (Itaipu), com melhorias recentes mas picos de instabilidade no verão.
\item[Água]  Abundante (Rio Paraguai).
\item[Solo]  Latossolos vermelhos férteis e profundos; excelente para silvicultura (eucalipto), gergelim e soja.
\item[Autossuficiência]  Elevada em proteína animal (indústria cárnea) e potencial agrícola de exportação.
\item[Segurança]  Zona de conflito contra a insurgência (EPP\footnote{\fineEPP}) no departamento. Taxa de homicídios elevada (18-20/\allowbreak 100k), mas concentrada em áreas rurais e narcotráfico. A área urbana possui forte presença militar (CODI/\allowbreak FTC).
\item[Leis Local: Capital departamental. Restrição de Fronteira]  Aplicação da Lei 2532/\allowbreak 05 (faixa de 50km).
\end{description}
\subsubsection{Indicadores Sociais}

\textbf{Fontes:} PNUD 2020, INE\footnote{\fineINE} Censo 2022, INE\footnote{\fineINE} EPHC 2023, Ministerio
Público 2024\\
\textbf{Nota:} Valores marcados como \emph{(dept.)} referem-se ao
departamento de Presidente Hayes; valores \emph{(dist.)} são específicos
deste distrito. Dados marcados \emph{(est.)} são estimativas.

{\def\LTcaptype{none} % do not increment counter
\begin{longtable}[]{@{}
  >{\raggedright\arraybackslash}p{(\linewidth - 4\tabcolsep) * \real{0.4231}}
  >{\raggedright\arraybackslash}p{(\linewidth - 4\tabcolsep) * \real{0.2692}}
  >{\raggedright\arraybackslash}p{(\linewidth - 4\tabcolsep) * \real{0.3077}}@{}}
\toprule\noalign{}
\begin{minipage}[b]{\linewidth}\raggedright
Indicador
\end{minipage} & \begin{minipage}[b]{\linewidth}\raggedright
Valor
\end{minipage} & \begin{minipage}[b]{\linewidth}\raggedright
Âmbito
\end{minipage} \\
\midrule\noalign{}
\endhead
\bottomrule\noalign{}
\endlastfoot
IDH (2020) & 0,693 & dept. \\
Ranking IDH nacional & 13/\allowbreak 18 & dept. \\
Esperança de vida & 71,8 anos & dept. \\
Escolaridade média & 8.8 anos & dist. \\
RNB per capita (USD PPA) & 7.900 & dept. \\
Pobreza monetária (\%) & 29,4\% & dept. \\
Pobreza extrema (\%) & 7,2\% & dept. \\
Índice de Gini & 0,453 & dept. \\
Acesso a água potável (\%) & 71,3\% & dept. \\
Acesso a saneamento (\%) & 68,9\% & dept. \\
Taxa de homicídios (est., /\allowbreak 100k hab) & 5,5--5,6 (2019--2020, fonte INE
--- próxima à média nacional) & dept. \\
Índice de segurança & médio & dept. \\
População (Censo 2022) & 47.967 habitantes & dist. \\
Idade mediana & N/\allowbreak D & dist. \\
Taxa de fecundidade & N/\allowbreak D & dist. \\
\end{longtable}
}
\subsubsection{Combustível}

\textbf{Referência:} PETROPAR /\allowbreak  postos locais (2024) \textbf{Tipo de
localidade:} Capital departamental

{\def\LTcaptype{none} % do not increment counter
\begin{longtable}[]{@{}lll@{}}
\toprule\noalign{}
Combustível & USD/\allowbreak litro & Gs/\allowbreak litro (aprox.) \\
\midrule\noalign{}
\endhead
\bottomrule\noalign{}
\endlastfoot
Gasolina 93 oct & 1.05 & 7,770 \\
Gasolina 97 oct (premium) & 1.15 & 8,510 \\
Diesel & 0.98 & 7,252 \\
\end{longtable}
}

\begin{quote}
Preços podem variar ±5\% conforme posto e sazonalidade. Chaco e interior
remoto apresentam maior variação.
\end{quote}

\subsubsection{Cobertura Celular}

\textbf{Fonte:} CONATEL PY /\allowbreak  operadoras (2024)

{\def\LTcaptype{none} % do not increment counter
\begin{longtable}[]{@{}ll@{}}
\toprule\noalign{}
Parâmetro & Valor \\
\midrule\noalign{}
\endhead
\bottomrule\noalign{}
\endlastfoot
Cobertura 4G população (dept.) & 60\% \\
Cobertura 4G área rural & 35\% \\
Melhor operadora & Tigo \\
Qualidade rural & limitada \\
\end{longtable}
}

\begin{quote}
Para áreas rurais fora do núcleo urbano, recomenda-se chip Tigo como
principal e Personal como backup.
\end{quote}

\subsubsection{Internet}

\textbf{Fonte:} CONATEL /\allowbreak  Speedtest Ookla (2024)

{\def\LTcaptype{none} % do not increment counter
\begin{longtable}[]{@{}ll@{}}
\toprule\noalign{}
Parâmetro & Valor \\
\midrule\noalign{}
\endhead
\bottomrule\noalign{}
\endlastfoot
Velocidade média download & 25 Mbps \\
Domicílios com internet (dept.) & 35\% \\
Tecnologia predominante & rádio/\allowbreak satélite \\
Opção rural & Starlink disponível (\textasciitilde USD 44/\allowbreak mês) \\
\end{longtable}
}

\subsubsection{Mercado Imobiliário e Terra
Rural}

\textbf{Fonte:} INDERT /\allowbreak  Clasificados.com.py (2024)

{\def\LTcaptype{none} % do not increment counter
\begin{longtable}[]{@{}ll@{}}
\toprule\noalign{}
Tipo & Referência \\
\midrule\noalign{}
\endhead
\bottomrule\noalign{}
\endlastfoot
Terra agrícola alta prod. (USD/\allowbreak ha) & 2,000 \\
Imóvel urbano (USD/\allowbreak m²) & 690 \\
Aluguel 2 quartos (USD/\allowbreak mês) & 300 \\
\end{longtable}
}

\begin{quote}
Valores de referência departamental. Localidades menores podem ter
preços 20--40\% abaixo da capital departamental.
\end{quote}

\subsubsection{Saúde}

\textbf{Fonte:} MSPBS /\allowbreak  IPS Paraguay (2024-2026), consolidação
departamental e proxy local

{\def\LTcaptype{none} % do not increment counter
\begin{longtable}[]{@{}ll@{}}
\toprule\noalign{}
Serviço & Disponibilidade \\
\midrule\noalign{}
\endhead
\bottomrule\noalign{}
\endlastfoot
USF /\allowbreak  Posto de Saúde & sim \\
Hospital Regional & sim \\
IPS (seguro social) & sim \\
Farmácia & sim \\
Distância ao hospital de referência & local \\
\end{longtable}
}

\textbf{Principais estabelecimentos:} Hospital distrital /\allowbreak  regional de
Villa Hayes; rede de USF e farmacias locais

\textbf{Observação para imigrantes:} Centro de referencia do
departamento. Acesso a atendimento primario e, em geral, a maior oferta
publica e privada da area. Cobertura privada continua recomendada para
especialidades e urgências de maior complexidade.
\newpage
\secaoDiagnostico{Horqueta}{\mapaDistrito{0103}}{Horqueta é a potência agrícola do departamento de Concepción. Oferece
alguns dos melhores solos do Paraguai e uma infraestrutura de serviços
superior à média regional. Seu GSS de 6.2 equilibra a excepcional
capacidade produtiva com os desafios de segurança tática inerentes à
zona de atuação da FTC.

\subsubsection{Por que sim}

\subsubsection{Por que não}

\subsubsection{Recomendação
Técnica}

Altamente indicado para investimentos em agronegócio e projetos de
autossuficiência que exijam solos de qualidade superior. Requer plano de
segurança patrimonial e integração com as autoridades locais (FTC/\allowbreak CODI).
Referência atual de score: GSS 6.2.

\begin{itemize}
\tightlist
\item
  \begin{enumerate}
  \def\labelenumi{\arabic{enumi})}
  \tightlist
  \item
    Dados censitários (INE\footnote{\fineINE} 2022).
  \end{enumerate}
\item
  \begin{enumerate}
  \def\labelenumi{\arabic{enumi})}
  \setcounter{enumi}{1}
  \tightlist
  \item
    Relatórios de produção agrícola (MAG/\allowbreak Indústria Stevia).
  \end{enumerate}
\item
  \begin{enumerate}
  \def\labelenumi{\arabic{enumi})}
  \setcounter{enumi}{2}
  \tightlist
  \item
    Mapas de segurança regional (CODI 2025).
  \end{enumerate}
\item
  \begin{enumerate}
  \def\labelenumi{\arabic{enumi})}
  \setcounter{enumi}{3}
  \tightlist
  \item
    Lei de Segurança Fronteiriça (2532/\allowbreak 05).
  \end{enumerate}
\end{itemize}}{dist:01_Concepción:Horqueta}
\begin{table}[ht!]
\small \begin{tabular}{p{3.0cm}p{0.8cm}p{6.0cm}} \toprule
\textbf{Critério} & \textbf{Nota} & \textbf{Justificativa} \\ \midrule
A - Ameaças Estratégicas & 5.0 & Hub agrícola estratégico e base militar tática regional \\
B - Risco Social & 5.5 & Densidade urbana equilibrada com presença constante de forças de segurança \\
C - Riscos Naturais & 7.0 & Geologia estável; relevo favorece drenagem e evita inundações graves \\
D - Autossuficiência & 8.0 & Solo de altíssima produtividade; abundância de recursos alimentares e hídricos \\
E - Institucional & 6.0 & Ambiente agroindustrial forte; restrições legais para estrangeiros brasileiros \\
\midrule \textbf{GSS FINAL} & \textbf{6.2} & \textbf{Moderadamente Seguro (Ideal para perfis agroindustriais com foco em produtividade e autonomia).} \\ \bottomrule \end{tabular}
\end{table}
\paragraph{Solo (SoilGrids 2.0, média ponderada 0--30
cm)}

{\def\LTcaptype{none} % do not increment counter
\begin{longtable}[]{@{}ll@{}}
\toprule\noalign{}
Parâmetro & Valor \\
\midrule\noalign{}
\endhead
\bottomrule\noalign{}
\endlastfoot
pH (H$_2$O) & 6.85 \\
Carbono orgânico (SOC) & 15.05 g/\allowbreak kg \\
Argila & 30.55 \% \\
Areia & 35.55 \% \\
Silte (calc.) & 33.9 \% \\
Densidade aparente & 1.4 g/\allowbreak cm³ \\
\textbf{Aptidão agrícola} & \textbf{Alta} \\
\end{longtable}
}

\subsubsection{Indicadores Sociais}

\textbf{Fontes:} PNUD 2020, INE\footnote{\fineINE} Censo 2022, INE\footnote{\fineINE} EPHC 2023, Ministerio
Público 2024\\
\textbf{Nota:} Valores marcados como \emph{(dept.)} referem-se ao
departamento de Boquerón; valores \emph{(dist.)} são específicos deste
distrito. Dados marcados \emph{(est.)} são estimativas.

{\def\LTcaptype{none} % do not increment counter
\begin{longtable}[]{@{}
  >{\raggedright\arraybackslash}p{(\linewidth - 4\tabcolsep) * \real{0.4231}}
  >{\raggedright\arraybackslash}p{(\linewidth - 4\tabcolsep) * \real{0.2692}}
  >{\raggedright\arraybackslash}p{(\linewidth - 4\tabcolsep) * \real{0.3077}}@{}}
\toprule\noalign{}
\begin{minipage}[b]{\linewidth}\raggedright
Indicador
\end{minipage} & \begin{minipage}[b]{\linewidth}\raggedright
Valor
\end{minipage} & \begin{minipage}[b]{\linewidth}\raggedright
Âmbito
\end{minipage} \\
\midrule\noalign{}
\endhead
\bottomrule\noalign{}
\endlastfoot
IDH (2020) & 0,676 & dept. \\
Ranking IDH nacional & 16/\allowbreak 18 & dept. \\
Esperança de vida & 70,8 anos & dept. \\
Escolaridade média & 7.2 anos & dist. \\
RNB per capita (USD PPA) & 8.600 & dept. \\
Pobreza monetária (\%) & estimativa 24,0\% & dept. \\
Pobreza extrema (\%) & estimativa 7,5\% & dept. \\
Índice de Gini & estimativa 0,485 & dept. \\
Acesso a água potável (\%) & estimativa 52,3\% & dept. \\
Acesso a saneamento (\%) & estimativa 48,7\% & dept. \\
Taxa de homicídios (est., /\allowbreak 100k hab) & 0,7--1,7 (2019--2020, fonte INE
--- entre as mais baixas do país) & dept. \\
Índice de segurança & alto & dept. \\
População (Censo 2022) & N/\allowbreak D & dist. \\
Idade mediana & N/\allowbreak D & dist. \\
Taxa de fecundidade & N/\allowbreak D & dist. \\
\end{longtable}
}
\subsubsection{Combustível}

\textbf{Referência:} PETROPAR /\allowbreak  postos locais (2024) \textbf{Tipo de
localidade:} Interior

{\def\LTcaptype{none} % do not increment counter
\begin{longtable}[]{@{}lll@{}}
\toprule\noalign{}
Combustível & USD/\allowbreak litro & Gs/\allowbreak litro (aprox.) \\
\midrule\noalign{}
\endhead
\bottomrule\noalign{}
\endlastfoot
Gasolina 93 oct & 1.12 & 8,288 \\
Gasolina 97 oct (premium) & 1.22 & 9,028 \\
Diesel & 1.05 & 7,770 \\
\end{longtable}
}

\begin{quote}
Preços podem variar ±5\% conforme posto e sazonalidade. Chaco e interior
remoto apresentam maior variação.
\end{quote}

\subsubsection{Cobertura Celular}

\textbf{Fonte:} CONATEL PY /\allowbreak  operadoras (2024)

{\def\LTcaptype{none} % do not increment counter
\begin{longtable}[]{@{}ll@{}}
\toprule\noalign{}
Parâmetro & Valor \\
\midrule\noalign{}
\endhead
\bottomrule\noalign{}
\endlastfoot
Cobertura 4G população (dept.) & 55\% \\
Cobertura 4G área rural & 30\% \\
Melhor operadora & Tigo \\
Qualidade rural & limitada \\
\end{longtable}
}

\begin{quote}
Para áreas rurais fora do núcleo urbano, recomenda-se chip Tigo como
principal e Personal como backup.
\end{quote}

\subsubsection{Internet}

\textbf{Fonte:} CONATEL /\allowbreak  Speedtest Ookla (2024)

{\def\LTcaptype{none} % do not increment counter
\begin{longtable}[]{@{}ll@{}}
\toprule\noalign{}
Parâmetro & Valor \\
\midrule\noalign{}
\endhead
\bottomrule\noalign{}
\endlastfoot
Velocidade média download & 28 Mbps \\
Domicílios com internet (dept.) & 38\% \\
Tecnologia predominante & rádio/\allowbreak satélite \\
Opção rural & Starlink disponível (\textasciitilde USD 44/\allowbreak mês) \\
\end{longtable}
}

\subsubsection{Mercado Imobiliário e Terra
Rural}

\textbf{Fonte:} INDERT /\allowbreak  Clasificados.com.py (2024)

{\def\LTcaptype{none} % do not increment counter
\begin{longtable}[]{@{}ll@{}}
\toprule\noalign{}
Tipo & Referência \\
\midrule\noalign{}
\endhead
\bottomrule\noalign{}
\endlastfoot
Terra agrícola alta prod. (USD/\allowbreak ha) & 1,200 \\
Imóvel urbano (USD/\allowbreak m²) & 700 \\
Aluguel 2 quartos (USD/\allowbreak mês) & 300 \\
\end{longtable}
}

\begin{quote}
Valores de referência departamental. Localidades menores podem ter
preços 20--40\% abaixo da capital departamental.
\end{quote}

\subsubsection{Saúde}

\textbf{Fonte:} MSPBS /\allowbreak  IPS Paraguay (2024-2026), consolidação
departamental e proxy local

{\def\LTcaptype{none} % do not increment counter
\begin{longtable}[]{@{}ll@{}}
\toprule\noalign{}
Serviço & Disponibilidade \\
\midrule\noalign{}
\endhead
\bottomrule\noalign{}
\endlastfoot
USF /\allowbreak  Posto de Saúde & sim \\
Hospital Regional & não \\
IPS (seguro social) & não \\
Farmácia & sim \\
Distância ao hospital de referência & \textasciitilde62 km
(Filadelfia) \\
\end{longtable}
}

\textbf{Principais estabelecimentos:} USF local; referencia hospitalar
em Filadelfia

\textbf{Observação para imigrantes:} Atendimento primario local ou em
raio curto; casos de maior complexidade seguem para Filadelfia.
Cobertura privada continua recomendada para especialidades e urgências
de maior complexidade.
\newpage
\secaoDiagnostico{Itacua}{\mapaDistrito{0114}}{Itacuá é uma localidade de fronteira mineral, oferecendo isolamento
geográfico massivo e abundância de água e minerais (calcário). Sua
viabilidade como refúgio é limitada pelo alto risco hidrológico e
restrições legais para estrangeiros.

\subsubsection{Por que sim}

\begin{itemize}
\tightlist
\item
  Densidade demográfica extremamente baixa (isolamento).
\item
  Recursos minerais abundantes e custo de vida reduzido.
\end{itemize}

\subsubsection{Por que não}

\subsubsection{Recomendação
Técnica}

Indicado como local de apoio ou exploração mineral específica. Para
moradia permanente, exige infraestrutura de defesa contra inundações e
autonomia logística total. Referência atual de score: GSS 5.8.

\begin{itemize}
\tightlist
\item
  \begin{enumerate}
  \def\labelenumi{\arabic{enumi})}
  \tightlist
  \item
    Censo INE\footnote{\fineINE} 2022.
  \end{enumerate}
\item
  \begin{enumerate}
  \def\labelenumi{\arabic{enumi})}
  \setcounter{enumi}{1}
  \tightlist
  \item
    Relatórios de cheias SEN\footnote{\fineSEN} 2019-2023.
  \end{enumerate}
\item
  \begin{enumerate}
  \def\labelenumi{\arabic{enumi})}
  \setcounter{enumi}{2}
  \tightlist
  \item
    Dados da indústria mineral regional (INC/\allowbreak Itakui).
  \end{enumerate}
\item
  \begin{enumerate}
  \def\labelenumi{\arabic{enumi})}
  \setcounter{enumi}{3}
  \tightlist
  \item
    Lei de Segurança Fronteiriça (2532/\allowbreak 05).
  \end{enumerate}
\end{itemize}}{dist:01_Concepción:Itacua}
\begin{table}[ht!]
\small \begin{tabular}{p{3.0cm}p{0.8cm}p{6.0cm}} \toprule
\textbf{Critério} & \textbf{Nota} & \textbf{Justificativa} \\ \midrule
A - Ameaças Estratégicas & 6.0 & Remoto, mas porto industrial e proximidade de fronteira fluvial são fatores táticos \\
B - Risco Social & 6.5 & Densidade baixíssima compensada pela instabilidade regional do departamento \\
C - Riscos Naturais & 5.5 & Risco crítico de inundação fluvial cíclica \\
D - Autossuficiência & 6.0 & Solo calcário dominante; abundância hídrica mas isolamento de recursos agrícolas variados \\
E - Institucional & 5.0 & Novo município; restrições de fronteira para estrangeiros limítrofes \\
\midrule \textbf{GSS FINAL} & \textbf{5.8} & \textbf{Moderadamente Seguro (Recomendado apenas para perfis com foco em mineração ou isolamento fluvial).} \\ \bottomrule \end{tabular}
\end{table}
\subsubsection{Vulnerabilidades e
Mitigação}

\begin{itemize}
\tightlist
\item
  \textbf{Isolamento Geográfico:} Dependência de transporte fluvial e
  estradas precárias (Mitigação: manter embarcação própria e estoque de
  suprimentos para 30 dias).
\item
  \textbf{Risco de Inundação:} Vulnerabilidade extrema a cheias do Rio
  Paraguai (Mitigação: construção em cotas elevadas, evitar áreas
  ribeirinhas diretas).
\item
  \textbf{Instabilidade Energética:} Final de linha de transmissão
  (Mitigação: sistemas solares de backup com baterias).
\item
  \textbf{Segurança Regional:} Monitoramento de atividades de grupos
  armados no departamento.
\end{itemize}

\subsubsection{Dossiê de Campo}
\begin{description}[style=nextline,leftmargin=0.5cm,font=\bfseries]
\item[Coordenadas]  22°59'S 57°43'W.
\item[Topografia]  Margem esquerda do Rio Paraguai. Relevo predominantemente plano com formações calcárias. Geologicamente estável, risco sísmico muito baixo.
\item[Alvos Estratégicos]  Indústrias de calcário (Calerías), Porto de Itacuá (escoamento de cimento e minerais), proximidade com a fronteira brasileira via rio.
\item[Fallout]  Ventos predominantemente do Norte/\allowbreak Nordeste (quentes/\allowbreak úmidos) e Sul (frentes frias).
\item[População]  2.377 (Censo 2022 Final).
\item[Densidade]  \textasciitilde 1,39 hab/\allowbreak km² (Extremamente baixa, excelente para isolamento estratégico).
\item[Vias de Acesso]  Acesso terrestre precário (Ruta PY22 e ramais de terra), forte dependência de transporte fluvial. Pista de pouso asfaltada em Vallemí (\textasciitilde 10 km).
\item[Serviços]  Infraestrutura básica local. Dependência de Concepción (190 km) ou Pedro Juan Caballero para saúde de alta complexidade.
\item[Custo de Vida]  \textasciitilde 30-40\% mais baixo que centros urbanos brasileiros; energia e água baratas.
\item[Preço da Terra]  US\$ 1.500-5.000/\allowbreak ha (rural/\allowbreak pastagem); lotes urbanos em San Lázaro por US\$ 15.000-30.000.
\item[Hidrologia]  Risco Crítico de Inundação. Proximidade direta com o Rio Paraguai; histórico de cheias severas (2019, 2023) que podem isolar a localidade por terra.
\item[Clima]  Tropical Úmido. Precipitação anual de 1.300-1.500 mm. Verões quentes com tempestades severas.
\item[Energia]  Rede nacional (ANDE\footnote{\fineANDE}/\allowbreak Itaipu), sujeita a instabilidades por ser final de linha no Norte.
\item[Água]  Abundante (Rio Paraguai).
\item[Qualidade do Solo]  Solo calcário (carbonato de cálcio/\allowbreak magnésio); ideal para correção de solos ácidos, exige manejo para agricultura intensiva diversificada.
\item[Recursos Locais]  Extração mineral (calcário/\allowbreak cimento), pesca, agricultura de subsistência e pecuária.
\item[Segurança]  Zona remota com baixa criminalidade urbana comum. Presença da FTC (Força de Tarefa Conjunta) na região devido à atividade insurgente (EPP\footnote{\fineEPP}) no departamento. Taxa de homicídios regional de 14-19/\allowbreak 100k (focada em áreas rurais/\allowbreak narcotráfico).
\item[Leis Local: Município de criação recente (2021). Restrição de Fronteira (Lei 2532/\allowbreak 05)]  Proibição de compra de terras rurais por estrangeiros de países limítrofes (faixa de 50km).
\end{description}
\newpage
\secaoDiagnostico{Loreto}{\mapaDistrito{0104}}{Loreto (Concepción) apresenta perfil operacional consolidado para
comparacao territorial, com leitura conjunta de infraestrutura, risco
hidroclimatico e capacidade de continuidade de servicos essenciais.

\subsubsection{Por que sim}

\subsubsection{Por que não}

\subsubsection{Pre-condicoes Minimas de
Decisao}

\begin{itemize}
\tightlist
\item
  Atualizar indicadores criticos em ciclo trimestral.
\item
  Confirmar trafegabilidade e contingencia hidrica/\allowbreak energetica local.
\item
  Validar checklist de resiliencia domiciliar/\allowbreak logistica antes de decisao
  final.
\end{itemize}

\subsubsection{Recomendação
Técnica}

Uso recomendado como candidato em analise multicriterio, condicionado ao
cumprimento das pre-condicoes acima. Referência atual de score: GSS 7.9;
risco predominante: baixo.

\begin{itemize}
\tightlist
\item
  \begin{enumerate}
  \def\labelenumi{\arabic{enumi})}
  \tightlist
  \item
    Delimitacao territorial validada por georreferencia institucional.
  \end{enumerate}
\item
  \begin{enumerate}
  \def\labelenumi{\arabic{enumi})}
  \setcounter{enumi}{1}
  \tightlist
  \item
    População municipal/\allowbreak distrital referenciada por base censitaria
    nacional.
  \end{enumerate}
\item
  \begin{enumerate}
  \def\labelenumi{\arabic{enumi})}
  \setcounter{enumi}{2}
  \tightlist
  \item
    Comparabilidade intra-departamental apoiada em indicadores
    distritais oficiais.
  \end{enumerate}
\item
  \begin{enumerate}
  \def\labelenumi{\arabic{enumi})}
  \setcounter{enumi}{3}
  \tightlist
  \item
    Estado macro de conectividade viaria ancorado em informacoes de
    rodovias.
  \end{enumerate}
\item
  \begin{enumerate}
  \def\labelenumi{\arabic{enumi})}
  \setcounter{enumi}{4}
  \tightlist
  \item
    Exposicao hidrometeorologica monitorada por avisos oficiais.
  \end{enumerate}
\item
  \begin{enumerate}
  \def\labelenumi{\arabic{enumi})}
  \setcounter{enumi}{5}
  \tightlist
  \item
    Historico de contexto meteorologico apoiado em publicacoes tecnicas.
  \end{enumerate}
\item
  \begin{enumerate}
  \def\labelenumi{\arabic{enumi})}
  \setcounter{enumi}{6}
  \tightlist
  \item
    Capacidade institucional de resposta a eventos apoiada em acoes da
    SEN\footnote{\fineSEN}.
  \end{enumerate}
\item
  \begin{enumerate}
  \def\labelenumi{\arabic{enumi})}
  \setcounter{enumi}{7}
  \tightlist
  \item
    Infraestrutura energetica analisada via operador nacional.
  \end{enumerate}
\item
  \begin{enumerate}
  \def\labelenumi{\arabic{enumi})}
  \setcounter{enumi}{8}
  \tightlist
  \item
    Referências de obras e logistica nacional verificadas no MOPC.
  \end{enumerate}
\item
  \begin{enumerate}
  \def\labelenumi{\arabic{enumi})}
  \setcounter{enumi}{9}
  \tightlist
  \item
    Contexto sociopolitico institucional referenciado por autoridade
    eleitoral.
  \end{enumerate}
\item
  \begin{enumerate}
  \def\labelenumi{\arabic{enumi})}
  \setcounter{enumi}{10}
  \tightlist
  \item
    Camada cartografica de apoio eleitoral/\allowbreak territorial consultada.
  \end{enumerate}
\item
  \begin{enumerate}
  \def\labelenumi{\arabic{enumi})}
  \setcounter{enumi}{11}
  \tightlist
  \item
    Dados publicos complementares para verificacao cruzada.
  \end{enumerate}
\end{itemize}}{dist:01_Concepción:Loreto}
\begin{table}[ht!]
\small \begin{tabular}{p{3.0cm}p{0.8cm}p{6.0cm}} \toprule
\textbf{Critério} & \textbf{Nota} & \textbf{Justificativa} \\ \midrule
A - Ameaças Estratégicas & 8.2 & - \\
B - Risco Social & 8.4 & - \\
C - Riscos Naturais & 6.5 & - \\
D - Autossuficiência & 8.0 & - \\
E - Institucional & 8.2 & - \\
\midrule \textbf{GSS FINAL} & \textbf{7.9} & \textbf{Seguro (bom potencial com preparacao).} \\ \bottomrule \end{tabular}
\end{table}
\newpage
\secaoDiagnostico{Paso Barreto}{\mapaDistrito{0111}}{Paso Barreto (Concepción) apresenta perfil operacional consolidado para
comparacao territorial, com leitura conjunta de infraestrutura, risco
hidroclimatico e capacidade de continuidade de servicos essenciais.

\subsubsection{Por que sim}

\begin{itemize}
\tightlist
\item
  Base institucional de referencia disponivel e rastreavel.
\item
  Estrutura comparativa (A-E/\allowbreak GSS) consistente para decisao relativa.
\item
  Plano de mitigacao acionavel ja definido no dossie.
\end{itemize}

\subsubsection{Por que não}

\subsubsection{Pre-condicoes Minimas de
Decisao}

\begin{itemize}
\tightlist
\item
  Atualizar indicadores criticos em ciclo trimestral.
\item
  Confirmar trafegabilidade e contingencia hidrica/\allowbreak energetica local.
\item
  Validar checklist de resiliencia domiciliar/\allowbreak logistica antes de decisao
  final.
\end{itemize}

\subsubsection{Recomendação
Técnica}

Uso recomendado como candidato em analise multicriterio, condicionado ao
cumprimento das pre-condicoes acima. Referência atual de score: GSS 7.0;
risco predominante: moderado.

\begin{itemize}
\tightlist
\item
  \begin{enumerate}
  \def\labelenumi{\arabic{enumi})}
  \tightlist
  \item
    Delimitacao territorial validada por georreferencia institucional.
  \end{enumerate}
\item
  \begin{enumerate}
  \def\labelenumi{\arabic{enumi})}
  \setcounter{enumi}{1}
  \tightlist
  \item
    População municipal/\allowbreak distrital referenciada por base censitaria
    nacional.
  \end{enumerate}
\item
  \begin{enumerate}
  \def\labelenumi{\arabic{enumi})}
  \setcounter{enumi}{2}
  \tightlist
  \item
    Comparabilidade intra-departamental apoiada em indicadores
    distritais oficiais.
  \end{enumerate}
\item
  \begin{enumerate}
  \def\labelenumi{\arabic{enumi})}
  \setcounter{enumi}{3}
  \tightlist
  \item
    Estado macro de conectividade viaria ancorado em informacoes de
    rodovias.
  \end{enumerate}
\item
  \begin{enumerate}
  \def\labelenumi{\arabic{enumi})}
  \setcounter{enumi}{4}
  \tightlist
  \item
    Exposicao hidrometeorologica monitorada por avisos oficiais.
  \end{enumerate}
\item
  \begin{enumerate}
  \def\labelenumi{\arabic{enumi})}
  \setcounter{enumi}{5}
  \tightlist
  \item
    Historico de contexto meteorologico apoiado em publicacoes tecnicas.
  \end{enumerate}
\item
  \begin{enumerate}
  \def\labelenumi{\arabic{enumi})}
  \setcounter{enumi}{6}
  \tightlist
  \item
    Capacidade institucional de resposta a eventos apoiada em acoes da
    SEN\footnote{\fineSEN}.
  \end{enumerate}
\item
  \begin{enumerate}
  \def\labelenumi{\arabic{enumi})}
  \setcounter{enumi}{7}
  \tightlist
  \item
    Infraestrutura energetica analisada via operador nacional.
  \end{enumerate}
\item
  \begin{enumerate}
  \def\labelenumi{\arabic{enumi})}
  \setcounter{enumi}{8}
  \tightlist
  \item
    Referências de obras e logistica nacional verificadas no MOPC.
  \end{enumerate}
\item
  \begin{enumerate}
  \def\labelenumi{\arabic{enumi})}
  \setcounter{enumi}{9}
  \tightlist
  \item
    Contexto sociopolitico institucional referenciado por autoridade
    eleitoral.
  \end{enumerate}
\item
  \begin{enumerate}
  \def\labelenumi{\arabic{enumi})}
  \setcounter{enumi}{10}
  \tightlist
  \item
    Camada cartografica de apoio eleitoral/\allowbreak territorial consultada.
  \end{enumerate}
\item
  \begin{enumerate}
  \def\labelenumi{\arabic{enumi})}
  \setcounter{enumi}{11}
  \tightlist
  \item
    Dados publicos complementares para verificacao cruzada.
  \end{enumerate}
\end{itemize}}{dist:01_Concepción:Paso_Barreto}
\begin{table}[ht!]
\small \begin{tabular}{p{3.0cm}p{0.8cm}p{6.0cm}} \toprule
\textbf{Critério} & \textbf{Nota} & \textbf{Justificativa} \\ \midrule
A - Ameaças Estratégicas & 7.2 & - \\
B - Risco Social & 7.2 & - \\
C - Riscos Naturais & 6.4 & - \\
D - Autossuficiência & 6.4 & - \\
E - Institucional & 7.6 & - \\
\midrule \textbf{GSS FINAL} & \textbf{7.0} & \textbf{Seguro (bom potencial com preparacao).} \\ \bottomrule \end{tabular}
\end{table}
\subsubsection{Vulnerabilidades e
Mitigação}

\begin{itemize}
\tightlist
\item
  Exposicao hidrometeorologica moderada (45-64/\allowbreak 100).
\item
  Conectividade viaria intermediaria (50-69/\allowbreak 100).
\item
  Estoque de contingencia para 14 dias (agua/\allowbreak alimentos/\allowbreak combustivel),
  priorizado quando risco hidro=45/\allowbreak 100.
\item
  Duplicar rotas/\allowbreak logistica quando conectividade viaria=51/\allowbreak 100 for
  \textless70; manter rota secundaria mapeada e testada trimestralmente.
\item
  Plano de energia resiliente com autonomia minima de 72h
  (geracao/\allowbreak backup), revisao semestral.
\end{itemize}

Sintese aplicada: - Dados textuais consolidados com base nas fontes
oficiais acima; proximos ciclos podem aprofundar indicadores numericos
especificos por distrito.
\subsubsection{Dossiê de Campo}
\begin{description}[style=nextline,leftmargin=0.5cm,font=\bfseries]
\item[Coordenadas]  23°03'S 57°07'W (geocodificado via Nominatim).
\end{description}
\textbf{Fonte climática:} NASA POWER Climatology API (período 2001-2020)

\textbf{Fonte luminosa:} estimativa world atlas

\paragraph{Irradiação Solar
(kWh/\allowbreak m²/\allowbreak dia)}

{\def\LTcaptype{none} % do not increment counter
\begin{longtable}[]{@{}
  >{\raggedright\arraybackslash}p{(\linewidth - 24\tabcolsep) * \real{0.0746}}
  >{\raggedright\arraybackslash}p{(\linewidth - 24\tabcolsep) * \real{0.0746}}
  >{\raggedright\arraybackslash}p{(\linewidth - 24\tabcolsep) * \real{0.0746}}
  >{\raggedright\arraybackslash}p{(\linewidth - 24\tabcolsep) * \real{0.0746}}
  >{\raggedright\arraybackslash}p{(\linewidth - 24\tabcolsep) * \real{0.0746}}
  >{\raggedright\arraybackslash}p{(\linewidth - 24\tabcolsep) * \real{0.0746}}
  >{\raggedright\arraybackslash}p{(\linewidth - 24\tabcolsep) * \real{0.0746}}
  >{\raggedright\arraybackslash}p{(\linewidth - 24\tabcolsep) * \real{0.0746}}
  >{\raggedright\arraybackslash}p{(\linewidth - 24\tabcolsep) * \real{0.0746}}
  >{\raggedright\arraybackslash}p{(\linewidth - 24\tabcolsep) * \real{0.0746}}
  >{\raggedright\arraybackslash}p{(\linewidth - 24\tabcolsep) * \real{0.0746}}
  >{\raggedright\arraybackslash}p{(\linewidth - 24\tabcolsep) * \real{0.0746}}
  >{\raggedright\arraybackslash}p{(\linewidth - 24\tabcolsep) * \real{0.1045}}@{}}
\toprule\noalign{}
\begin{minipage}[b]{\linewidth}\raggedright
Jan
\end{minipage} & \begin{minipage}[b]{\linewidth}\raggedright
Fev
\end{minipage} & \begin{minipage}[b]{\linewidth}\raggedright
Mar
\end{minipage} & \begin{minipage}[b]{\linewidth}\raggedright
Abr
\end{minipage} & \begin{minipage}[b]{\linewidth}\raggedright
Mai
\end{minipage} & \begin{minipage}[b]{\linewidth}\raggedright
Jun
\end{minipage} & \begin{minipage}[b]{\linewidth}\raggedright
Jul
\end{minipage} & \begin{minipage}[b]{\linewidth}\raggedright
Ago
\end{minipage} & \begin{minipage}[b]{\linewidth}\raggedright
Set
\end{minipage} & \begin{minipage}[b]{\linewidth}\raggedright
Out
\end{minipage} & \begin{minipage}[b]{\linewidth}\raggedright
Nov
\end{minipage} & \begin{minipage}[b]{\linewidth}\raggedright
Dez
\end{minipage} & \begin{minipage}[b]{\linewidth}\raggedright
Média
\end{minipage} \\
\midrule\noalign{}
\endhead
\bottomrule\noalign{}
\endlastfoot
6.68 & 6.07 & 5.58 & 4.63 & 3.47 & 3.00 & 3.30 & 4.09 & 4.68 & 5.49 &
6.34 & 6.61 & \textbf{5.0} \\
\end{longtable}
}

\textbf{Inclinação solar recomendada:} 23° N (anual) · 33° N (inverno
jun-ago) · 13° N (verão nov-jan)

\paragraph{Precipitação (mm/\allowbreak mês)}

{\def\LTcaptype{none} % do not increment counter
\begin{longtable}[]{@{}
  >{\raggedright\arraybackslash}p{(\linewidth - 24\tabcolsep) * \real{0.0704}}
  >{\raggedright\arraybackslash}p{(\linewidth - 24\tabcolsep) * \real{0.0704}}
  >{\raggedright\arraybackslash}p{(\linewidth - 24\tabcolsep) * \real{0.0704}}
  >{\raggedright\arraybackslash}p{(\linewidth - 24\tabcolsep) * \real{0.0704}}
  >{\raggedright\arraybackslash}p{(\linewidth - 24\tabcolsep) * \real{0.0704}}
  >{\raggedright\arraybackslash}p{(\linewidth - 24\tabcolsep) * \real{0.0704}}
  >{\raggedright\arraybackslash}p{(\linewidth - 24\tabcolsep) * \real{0.0704}}
  >{\raggedright\arraybackslash}p{(\linewidth - 24\tabcolsep) * \real{0.0704}}
  >{\raggedright\arraybackslash}p{(\linewidth - 24\tabcolsep) * \real{0.0704}}
  >{\raggedright\arraybackslash}p{(\linewidth - 24\tabcolsep) * \real{0.0704}}
  >{\raggedright\arraybackslash}p{(\linewidth - 24\tabcolsep) * \real{0.0704}}
  >{\raggedright\arraybackslash}p{(\linewidth - 24\tabcolsep) * \real{0.0704}}
  >{\raggedright\arraybackslash}p{(\linewidth - 24\tabcolsep) * \real{0.1549}}@{}}
\toprule\noalign{}
\begin{minipage}[b]{\linewidth}\raggedright
Jan
\end{minipage} & \begin{minipage}[b]{\linewidth}\raggedright
Fev
\end{minipage} & \begin{minipage}[b]{\linewidth}\raggedright
Mar
\end{minipage} & \begin{minipage}[b]{\linewidth}\raggedright
Abr
\end{minipage} & \begin{minipage}[b]{\linewidth}\raggedright
Mai
\end{minipage} & \begin{minipage}[b]{\linewidth}\raggedright
Jun
\end{minipage} & \begin{minipage}[b]{\linewidth}\raggedright
Jul
\end{minipage} & \begin{minipage}[b]{\linewidth}\raggedright
Ago
\end{minipage} & \begin{minipage}[b]{\linewidth}\raggedright
Set
\end{minipage} & \begin{minipage}[b]{\linewidth}\raggedright
Out
\end{minipage} & \begin{minipage}[b]{\linewidth}\raggedright
Nov
\end{minipage} & \begin{minipage}[b]{\linewidth}\raggedright
Dez
\end{minipage} & \begin{minipage}[b]{\linewidth}\raggedright
Total/\allowbreak ano
\end{minipage} \\
\midrule\noalign{}
\endhead
\bottomrule\noalign{}
\endlastfoot
165 & 157 & 117 & 134 & 112 & 71 & 46 & 32 & 73 & 146 & 205 & 165 &
\textbf{1428 mm} \\
\end{longtable}
}

\paragraph{Poluição Luminosa}

{\def\LTcaptype{none} % do not increment counter
\begin{longtable}[]{@{}ll@{}}
\toprule\noalign{}
Parâmetro & Valor \\
\midrule\noalign{}
\endhead
\bottomrule\noalign{}
\endlastfoot
Escala Bortle & 3 --- Céu rural \\
Radiância artificial & 1.5 nW/\allowbreak cm²/\allowbreak sr \\
\end{longtable}
}
\paragraph{Solo (SoilGrids 2.0, média ponderada 0--30
cm)}

{\def\LTcaptype{none} % do not increment counter
\begin{longtable}[]{@{}ll@{}}
\toprule\noalign{}
Parâmetro & Valor \\
\midrule\noalign{}
\endhead
\bottomrule\noalign{}
\endlastfoot
pH (H$_2$O) & 5.75 \\
Carbono orgânico (SOC) & 15.32 g/\allowbreak kg \\
Argila & 21.42 \% \\
Areia & 54.93 \% \\
Silte (calc.) & 23.6 \% \\
Densidade aparente & 1.31 g/\allowbreak cm³ \\
\textbf{Aptidão agrícola} & \textbf{Alta} \\
\end{longtable}
}

\subsubsection{Indicadores Sociais}

\textbf{Fontes:} PNUD 2020, INE\footnote{\fineINE} Censo 2022, INE\footnote{\fineINE} EPHC 2023, Ministerio
Público 2024\\
\textbf{Nota:} Valores marcados como \emph{(dept.)} referem-se ao
departamento de Concepción; valores \emph{(dist.)} são específicos deste
distrito. Dados marcados \emph{(est.)} são estimativas.

{\def\LTcaptype{none} % do not increment counter
\begin{longtable}[]{@{}
  >{\raggedright\arraybackslash}p{(\linewidth - 4\tabcolsep) * \real{0.4231}}
  >{\raggedright\arraybackslash}p{(\linewidth - 4\tabcolsep) * \real{0.2692}}
  >{\raggedright\arraybackslash}p{(\linewidth - 4\tabcolsep) * \real{0.3077}}@{}}
\toprule\noalign{}
\begin{minipage}[b]{\linewidth}\raggedright
Indicador
\end{minipage} & \begin{minipage}[b]{\linewidth}\raggedright
Valor
\end{minipage} & \begin{minipage}[b]{\linewidth}\raggedright
Âmbito
\end{minipage} \\
\midrule\noalign{}
\endhead
\bottomrule\noalign{}
\endlastfoot
IDH (2020) & 0.635 & dist. (est.) \\
Ranking IDH nacional & 9/\allowbreak 18 & dept. \\
Esperança de vida & 72,4 anos & dept. \\
Escolaridade média & 7.2 anos & dist. \\
RNB per capita (USD PPA) & 8.400 & dept. \\
Pobreza monetária (\%) & 31,1\% & dept. \\
Pobreza extrema (\%) & 7,8\% & dept. \\
Índice de Gini & 0,460 & dept. \\
Acesso a água potável (\%) & 74,2\% & dept. \\
Acesso a saneamento (\%) & 71,8\% & dept. \\
Taxa de homicídios (est., /\allowbreak 100k hab) & \textasciitilde12--15
(estimativa, acima da média nacional) & dept. \\
Índice de segurança & baixo & dept. \\
População (Censo 2022) & 3.988 hab. & dist. \\
Idade mediana & 25.0 anos & dist. \\
Taxa de fecundidade & 2 & dist. \\
\end{longtable}
}
\subsubsection{Combustível}

\textbf{Referência:} PETROPAR /\allowbreak  postos locais (2024) \textbf{Tipo de
localidade:} Interior

{\def\LTcaptype{none} % do not increment counter
\begin{longtable}[]{@{}lll@{}}
\toprule\noalign{}
Combustível & USD/\allowbreak litro & Gs/\allowbreak litro (aprox.) \\
\midrule\noalign{}
\endhead
\bottomrule\noalign{}
\endlastfoot
Gasolina 93 oct & 1.01 & 7,474 \\
Gasolina 97 oct (premium) & 1.11 & 8,214 \\
Diesel & 0.94 & 6,956 \\
\end{longtable}
}

\begin{quote}
Preços podem variar ±5\% conforme posto e sazonalidade. Chaco e interior
remoto apresentam maior variação.
\end{quote}

\subsubsection{Cobertura Celular}

\textbf{Fonte:} CONATEL PY /\allowbreak  operadoras (2024)

{\def\LTcaptype{none} % do not increment counter
\begin{longtable}[]{@{}ll@{}}
\toprule\noalign{}
Parâmetro & Valor \\
\midrule\noalign{}
\endhead
\bottomrule\noalign{}
\endlastfoot
Cobertura 4G população (dept.) & 78\% \\
Cobertura 4G área rural & 55\% \\
Melhor operadora & Tigo \\
Qualidade rural & moderada \\
\end{longtable}
}

\begin{quote}
Para áreas rurais fora do núcleo urbano, recomenda-se chip Tigo como
principal e Personal como backup.
\end{quote}

\subsubsection{Internet}

\textbf{Fonte:} CONATEL /\allowbreak  Speedtest Ookla (2024)

{\def\LTcaptype{none} % do not increment counter
\begin{longtable}[]{@{}ll@{}}
\toprule\noalign{}
Parâmetro & Valor \\
\midrule\noalign{}
\endhead
\bottomrule\noalign{}
\endlastfoot
Velocidade média download & 38 Mbps \\
Domicílios com internet (dept.) & 48\% \\
Tecnologia predominante & rádio \\
Opção rural & Starlink disponível (\textasciitilde USD 44/\allowbreak mês) \\
\end{longtable}
}

\subsubsection{Mercado Imobiliário e Terra
Rural}

\textbf{Fonte:} INDERT /\allowbreak  Clasificados.com.py (2024)

{\def\LTcaptype{none} % do not increment counter
\begin{longtable}[]{@{}ll@{}}
\toprule\noalign{}
Tipo & Referência \\
\midrule\noalign{}
\endhead
\bottomrule\noalign{}
\endlastfoot
Terra agrícola alta prod. (USD/\allowbreak ha) & 3,500 \\
Imóvel urbano (USD/\allowbreak m²) & 650 \\
Aluguel 2 quartos (USD/\allowbreak mês) & 260 \\
\end{longtable}
}

\begin{quote}
Valores de referência departamental. Localidades menores podem ter
preços 20--40\% abaixo da capital departamental.
\end{quote}

\subsubsection{Saúde}

\textbf{Fonte:} MSPBS /\allowbreak  IPS Paraguay (2024-2026), consolidação
departamental e proxy local

{\def\LTcaptype{none} % do not increment counter
\begin{longtable}[]{@{}ll@{}}
\toprule\noalign{}
Serviço & Disponibilidade \\
\midrule\noalign{}
\endhead
\bottomrule\noalign{}
\endlastfoot
USF /\allowbreak  Posto de Saúde & sim \\
Hospital Regional & não \\
IPS (seguro social) & sim \\
Farmácia & sim \\
Distância ao hospital de referência & \textasciitilde65 km
(Concepción) \\
\end{longtable}
}

\textbf{Principais estabelecimentos:} USF local; referencia hospitalar
em Concepción

\textbf{Observação para imigrantes:} Atendimento primario local ou em
raio curto; casos de maior complexidade seguem para Concepción.
Cobertura privada continua recomendada para especialidades e urgências
de maior complexidade.
\newpage
\secaoDiagnostico{Paso Horqueta}{\mapaDistrito{0113}}{Paso Horqueta (Concepción) apresenta perfil operacional consolidado para
comparacao territorial, com leitura conjunta de infraestrutura, risco
hidroclimatico e capacidade de continuidade de servicos essenciais.

\subsubsection{Por que sim}

\subsubsection{Por que não}

\subsubsection{Pre-condicoes Minimas de
Decisao}

\begin{itemize}
\tightlist
\item
  Atualizar indicadores criticos em ciclo trimestral.
\item
  Confirmar trafegabilidade e contingencia hidrica/\allowbreak energetica local.
\item
  Validar checklist de resiliencia domiciliar/\allowbreak logistica antes de decisao
  final.
\end{itemize}

\subsubsection{Recomendação
Técnica}

Uso recomendado como candidato em analise multicriterio, condicionado ao
cumprimento das pre-condicoes acima. Referência atual de score: GSS 8.0;
risco predominante: baixo.

\begin{itemize}
\tightlist
\item
  \begin{enumerate}
  \def\labelenumi{\arabic{enumi})}
  \tightlist
  \item
    Delimitacao territorial validada por georreferencia institucional.
  \end{enumerate}
\item
  \begin{enumerate}
  \def\labelenumi{\arabic{enumi})}
  \setcounter{enumi}{1}
  \tightlist
  \item
    População municipal/\allowbreak distrital referenciada por base censitaria
    nacional.
  \end{enumerate}
\item
  \begin{enumerate}
  \def\labelenumi{\arabic{enumi})}
  \setcounter{enumi}{2}
  \tightlist
  \item
    Comparabilidade intra-departamental apoiada em indicadores
    distritais oficiais.
  \end{enumerate}
\item
  \begin{enumerate}
  \def\labelenumi{\arabic{enumi})}
  \setcounter{enumi}{3}
  \tightlist
  \item
    Estado macro de conectividade viaria ancorado em informacoes de
    rodovias.
  \end{enumerate}
\item
  \begin{enumerate}
  \def\labelenumi{\arabic{enumi})}
  \setcounter{enumi}{4}
  \tightlist
  \item
    Exposicao hidrometeorologica monitorada por avisos oficiais.
  \end{enumerate}
\item
  \begin{enumerate}
  \def\labelenumi{\arabic{enumi})}
  \setcounter{enumi}{5}
  \tightlist
  \item
    Historico de contexto meteorologico apoiado em publicacoes tecnicas.
  \end{enumerate}
\item
  \begin{enumerate}
  \def\labelenumi{\arabic{enumi})}
  \setcounter{enumi}{6}
  \tightlist
  \item
    Capacidade institucional de resposta a eventos apoiada em acoes da
    SEN\footnote{\fineSEN}.
  \end{enumerate}
\item
  \begin{enumerate}
  \def\labelenumi{\arabic{enumi})}
  \setcounter{enumi}{7}
  \tightlist
  \item
    Infraestrutura energetica analisada via operador nacional.
  \end{enumerate}
\item
  \begin{enumerate}
  \def\labelenumi{\arabic{enumi})}
  \setcounter{enumi}{8}
  \tightlist
  \item
    Referências de obras e logistica nacional verificadas no MOPC.
  \end{enumerate}
\item
  \begin{enumerate}
  \def\labelenumi{\arabic{enumi})}
  \setcounter{enumi}{9}
  \tightlist
  \item
    Contexto sociopolitico institucional referenciado por autoridade
    eleitoral.
  \end{enumerate}
\item
  \begin{enumerate}
  \def\labelenumi{\arabic{enumi})}
  \setcounter{enumi}{10}
  \tightlist
  \item
    Camada cartografica de apoio eleitoral/\allowbreak territorial consultada.
  \end{enumerate}
\item
  \begin{enumerate}
  \def\labelenumi{\arabic{enumi})}
  \setcounter{enumi}{11}
  \tightlist
  \item
    Dados publicos complementares para verificacao cruzada.
  \end{enumerate}
\end{itemize}}{dist:01_Concepción:Paso_Horqueta}
\begin{table}[ht!]
\small \begin{tabular}{p{3.0cm}p{0.8cm}p{6.0cm}} \toprule
\textbf{Critério} & \textbf{Nota} & \textbf{Justificativa} \\ \midrule
A - Ameaças Estratégicas & 8.3 & - \\
B - Risco Social & 8.7 & - \\
C - Riscos Naturais & 7.5 & - \\
D - Autossuficiência & 6.6 & - \\
E - Institucional & 8.7 & - \\
\midrule \textbf{GSS FINAL} & \textbf{8.0} & \textbf{Seguro (bom potencial com preparacao).} \\ \bottomrule \end{tabular}
\end{table}
\subsubsection{Vulnerabilidades e
Mitigação}

\begin{itemize}
\tightlist
\item
  Exposicao hidrometeorologica controlada (\textless45/\allowbreak 100).
\item
  Conectividade viaria favoravel (\textgreater=70/\allowbreak 100).
\item
  Estoque de contingencia para 14 dias (agua/\allowbreak alimentos/\allowbreak combustivel),
  priorizado quando risco hidro=31/\allowbreak 100.
\item
  Duplicar rotas/\allowbreak logistica quando conectividade viaria=73/\allowbreak 100 for
  \textless70; manter rota secundaria mapeada e testada trimestralmente.
\item
  Plano de energia resiliente com autonomia minima de 72h
  (geracao/\allowbreak backup), revisao semestral.
\end{itemize}

Sintese aplicada: - Dados textuais consolidados com base nas fontes
oficiais acima; proximos ciclos podem aprofundar indicadores numericos
especificos por distrito.
\subsubsection{Dossiê de Campo}
\begin{description}[style=nextline,leftmargin=0.5cm,font=\bfseries]
\item[Coordenadas]  23°05'S 57°23'W (geocodificado via Nominatim).
\end{description}
\textbf{Fonte climática:} NASA POWER Climatology API (período 2001-2020)

\textbf{Fonte luminosa:} estimativa world atlas

\paragraph{Irradiação Solar
(kWh/\allowbreak m²/\allowbreak dia)}

{\def\LTcaptype{none} % do not increment counter
\begin{longtable}[]{@{}
  >{\raggedright\arraybackslash}p{(\linewidth - 24\tabcolsep) * \real{0.0746}}
  >{\raggedright\arraybackslash}p{(\linewidth - 24\tabcolsep) * \real{0.0746}}
  >{\raggedright\arraybackslash}p{(\linewidth - 24\tabcolsep) * \real{0.0746}}
  >{\raggedright\arraybackslash}p{(\linewidth - 24\tabcolsep) * \real{0.0746}}
  >{\raggedright\arraybackslash}p{(\linewidth - 24\tabcolsep) * \real{0.0746}}
  >{\raggedright\arraybackslash}p{(\linewidth - 24\tabcolsep) * \real{0.0746}}
  >{\raggedright\arraybackslash}p{(\linewidth - 24\tabcolsep) * \real{0.0746}}
  >{\raggedright\arraybackslash}p{(\linewidth - 24\tabcolsep) * \real{0.0746}}
  >{\raggedright\arraybackslash}p{(\linewidth - 24\tabcolsep) * \real{0.0746}}
  >{\raggedright\arraybackslash}p{(\linewidth - 24\tabcolsep) * \real{0.0746}}
  >{\raggedright\arraybackslash}p{(\linewidth - 24\tabcolsep) * \real{0.0746}}
  >{\raggedright\arraybackslash}p{(\linewidth - 24\tabcolsep) * \real{0.0746}}
  >{\raggedright\arraybackslash}p{(\linewidth - 24\tabcolsep) * \real{0.1045}}@{}}
\toprule\noalign{}
\begin{minipage}[b]{\linewidth}\raggedright
Jan
\end{minipage} & \begin{minipage}[b]{\linewidth}\raggedright
Fev
\end{minipage} & \begin{minipage}[b]{\linewidth}\raggedright
Mar
\end{minipage} & \begin{minipage}[b]{\linewidth}\raggedright
Abr
\end{minipage} & \begin{minipage}[b]{\linewidth}\raggedright
Mai
\end{minipage} & \begin{minipage}[b]{\linewidth}\raggedright
Jun
\end{minipage} & \begin{minipage}[b]{\linewidth}\raggedright
Jul
\end{minipage} & \begin{minipage}[b]{\linewidth}\raggedright
Ago
\end{minipage} & \begin{minipage}[b]{\linewidth}\raggedright
Set
\end{minipage} & \begin{minipage}[b]{\linewidth}\raggedright
Out
\end{minipage} & \begin{minipage}[b]{\linewidth}\raggedright
Nov
\end{minipage} & \begin{minipage}[b]{\linewidth}\raggedright
Dez
\end{minipage} & \begin{minipage}[b]{\linewidth}\raggedright
Média
\end{minipage} \\
\midrule\noalign{}
\endhead
\bottomrule\noalign{}
\endlastfoot
6.68 & 6.07 & 5.58 & 4.63 & 3.47 & 3.00 & 3.30 & 4.09 & 4.68 & 5.49 &
6.34 & 6.61 & \textbf{5.0} \\
\end{longtable}
}

\textbf{Inclinação solar recomendada:} 23° N (anual) · 33° N (inverno
jun-ago) · 13° N (verão nov-jan)

\paragraph{Precipitação (mm/\allowbreak mês)}

{\def\LTcaptype{none} % do not increment counter
\begin{longtable}[]{@{}
  >{\raggedright\arraybackslash}p{(\linewidth - 24\tabcolsep) * \real{0.0704}}
  >{\raggedright\arraybackslash}p{(\linewidth - 24\tabcolsep) * \real{0.0704}}
  >{\raggedright\arraybackslash}p{(\linewidth - 24\tabcolsep) * \real{0.0704}}
  >{\raggedright\arraybackslash}p{(\linewidth - 24\tabcolsep) * \real{0.0704}}
  >{\raggedright\arraybackslash}p{(\linewidth - 24\tabcolsep) * \real{0.0704}}
  >{\raggedright\arraybackslash}p{(\linewidth - 24\tabcolsep) * \real{0.0704}}
  >{\raggedright\arraybackslash}p{(\linewidth - 24\tabcolsep) * \real{0.0704}}
  >{\raggedright\arraybackslash}p{(\linewidth - 24\tabcolsep) * \real{0.0704}}
  >{\raggedright\arraybackslash}p{(\linewidth - 24\tabcolsep) * \real{0.0704}}
  >{\raggedright\arraybackslash}p{(\linewidth - 24\tabcolsep) * \real{0.0704}}
  >{\raggedright\arraybackslash}p{(\linewidth - 24\tabcolsep) * \real{0.0704}}
  >{\raggedright\arraybackslash}p{(\linewidth - 24\tabcolsep) * \real{0.0704}}
  >{\raggedright\arraybackslash}p{(\linewidth - 24\tabcolsep) * \real{0.1549}}@{}}
\toprule\noalign{}
\begin{minipage}[b]{\linewidth}\raggedright
Jan
\end{minipage} & \begin{minipage}[b]{\linewidth}\raggedright
Fev
\end{minipage} & \begin{minipage}[b]{\linewidth}\raggedright
Mar
\end{minipage} & \begin{minipage}[b]{\linewidth}\raggedright
Abr
\end{minipage} & \begin{minipage}[b]{\linewidth}\raggedright
Mai
\end{minipage} & \begin{minipage}[b]{\linewidth}\raggedright
Jun
\end{minipage} & \begin{minipage}[b]{\linewidth}\raggedright
Jul
\end{minipage} & \begin{minipage}[b]{\linewidth}\raggedright
Ago
\end{minipage} & \begin{minipage}[b]{\linewidth}\raggedright
Set
\end{minipage} & \begin{minipage}[b]{\linewidth}\raggedright
Out
\end{minipage} & \begin{minipage}[b]{\linewidth}\raggedright
Nov
\end{minipage} & \begin{minipage}[b]{\linewidth}\raggedright
Dez
\end{minipage} & \begin{minipage}[b]{\linewidth}\raggedright
Total/\allowbreak ano
\end{minipage} \\
\midrule\noalign{}
\endhead
\bottomrule\noalign{}
\endlastfoot
164 & 151 & 121 & 135 & 108 & 67 & 43 & 29 & 68 & 141 & 209 & 165 &
\textbf{1405 mm} \\
\end{longtable}
}

\paragraph{Poluição Luminosa}

{\def\LTcaptype{none} % do not increment counter
\begin{longtable}[]{@{}ll@{}}
\toprule\noalign{}
Parâmetro & Valor \\
\midrule\noalign{}
\endhead
\bottomrule\noalign{}
\endlastfoot
Escala Bortle & 3 --- Céu rural \\
Radiância artificial & 1.5 nW/\allowbreak cm²/\allowbreak sr \\
\end{longtable}
}
\paragraph{Solo (SoilGrids 2.0, média ponderada 0--30
cm)}

{\def\LTcaptype{none} % do not increment counter
\begin{longtable}[]{@{}ll@{}}
\toprule\noalign{}
Parâmetro & Valor \\
\midrule\noalign{}
\endhead
\bottomrule\noalign{}
\endlastfoot
pH (H$_2$O) & 5.88 \\
Carbono orgânico (SOC) & 16.02 g/\allowbreak kg \\
Argila & 19.35 \% \\
Areia & 53.9 \% \\
Silte (calc.) & 26.8 \% \\
Densidade aparente & 1.33 g/\allowbreak cm³ \\
\textbf{Aptidão agrícola} & \textbf{Alta} \\
\end{longtable}
}

\subsubsection{Indicadores Sociais}

\textbf{Fontes:} PNUD 2020, INE\footnote{\fineINE} Censo 2022, INE\footnote{\fineINE} EPHC 2023, Ministerio
Público 2024\\
\textbf{Nota:} Valores marcados como \emph{(dept.)} referem-se ao
departamento de Concepción; valores \emph{(dist.)} são específicos deste
distrito. Dados marcados \emph{(est.)} são estimativas.

{\def\LTcaptype{none} % do not increment counter
\begin{longtable}[]{@{}
  >{\raggedright\arraybackslash}p{(\linewidth - 4\tabcolsep) * \real{0.4231}}
  >{\raggedright\arraybackslash}p{(\linewidth - 4\tabcolsep) * \real{0.2692}}
  >{\raggedright\arraybackslash}p{(\linewidth - 4\tabcolsep) * \real{0.3077}}@{}}
\toprule\noalign{}
\begin{minipage}[b]{\linewidth}\raggedright
Indicador
\end{minipage} & \begin{minipage}[b]{\linewidth}\raggedright
Valor
\end{minipage} & \begin{minipage}[b]{\linewidth}\raggedright
Âmbito
\end{minipage} \\
\midrule\noalign{}
\endhead
\bottomrule\noalign{}
\endlastfoot
IDH (2020) & 0.640 & dist. (est.) \\
Ranking IDH nacional & 9/\allowbreak 18 & dept. \\
Esperança de vida & 72,4 anos & dept. \\
Escolaridade média & 7.4 anos & dist. \\
RNB per capita (USD PPA) & 8.400 & dept. \\
Pobreza monetária (\%) & 31,1\% & dept. \\
Pobreza extrema (\%) & 7,8\% & dept. \\
Índice de Gini & 0,460 & dept. \\
Acesso a água potável (\%) & 74,2\% & dept. \\
Acesso a saneamento (\%) & 71,8\% & dept. \\
Taxa de homicídios (est., /\allowbreak 100k hab) & \textasciitilde12--15
(estimativa, acima da média nacional) & dept. \\
Índice de segurança & baixo & dept. \\
População (Censo 2022) & 7.646 hab. & dist. \\
Idade mediana & 27.0 anos & dist. \\
Taxa de fecundidade & 2 & dist. \\
\end{longtable}
}
\subsubsection{Combustível}

\textbf{Referência:} PETROPAR /\allowbreak  postos locais (2024) \textbf{Tipo de
localidade:} Interior

{\def\LTcaptype{none} % do not increment counter
\begin{longtable}[]{@{}lll@{}}
\toprule\noalign{}
Combustível & USD/\allowbreak litro & Gs/\allowbreak litro (aprox.) \\
\midrule\noalign{}
\endhead
\bottomrule\noalign{}
\endlastfoot
Gasolina 93 oct & 1.01 & 7,474 \\
Gasolina 97 oct (premium) & 1.11 & 8,214 \\
Diesel & 0.94 & 6,956 \\
\end{longtable}
}

\begin{quote}
Preços podem variar ±5\% conforme posto e sazonalidade. Chaco e interior
remoto apresentam maior variação.
\end{quote}

\subsubsection{Cobertura Celular}

\textbf{Fonte:} CONATEL PY /\allowbreak  operadoras (2024)

{\def\LTcaptype{none} % do not increment counter
\begin{longtable}[]{@{}ll@{}}
\toprule\noalign{}
Parâmetro & Valor \\
\midrule\noalign{}
\endhead
\bottomrule\noalign{}
\endlastfoot
Cobertura 4G população (dept.) & 78\% \\
Cobertura 4G área rural & 55\% \\
Melhor operadora & Tigo \\
Qualidade rural & moderada \\
\end{longtable}
}

\begin{quote}
Para áreas rurais fora do núcleo urbano, recomenda-se chip Tigo como
principal e Personal como backup.
\end{quote}

\subsubsection{Internet}

\textbf{Fonte:} CONATEL /\allowbreak  Speedtest Ookla (2024)

{\def\LTcaptype{none} % do not increment counter
\begin{longtable}[]{@{}ll@{}}
\toprule\noalign{}
Parâmetro & Valor \\
\midrule\noalign{}
\endhead
\bottomrule\noalign{}
\endlastfoot
Velocidade média download & 38 Mbps \\
Domicílios com internet (dept.) & 48\% \\
Tecnologia predominante & rádio \\
Opção rural & Starlink disponível (\textasciitilde USD 44/\allowbreak mês) \\
\end{longtable}
}

\subsubsection{Mercado Imobiliário e Terra
Rural}

\textbf{Fonte:} INDERT /\allowbreak  Clasificados.com.py (2024)

{\def\LTcaptype{none} % do not increment counter
\begin{longtable}[]{@{}ll@{}}
\toprule\noalign{}
Tipo & Referência \\
\midrule\noalign{}
\endhead
\bottomrule\noalign{}
\endlastfoot
Terra agrícola alta prod. (USD/\allowbreak ha) & 3,500 \\
Imóvel urbano (USD/\allowbreak m²) & 650 \\
Aluguel 2 quartos (USD/\allowbreak mês) & 260 \\
\end{longtable}
}

\begin{quote}
Valores de referência departamental. Localidades menores podem ter
preços 20--40\% abaixo da capital departamental.
\end{quote}

\subsubsection{Saúde}

\textbf{Fonte:} MSPBS /\allowbreak  IPS Paraguay (2024-2026), consolidação
departamental e proxy local

{\def\LTcaptype{none} % do not increment counter
\begin{longtable}[]{@{}ll@{}}
\toprule\noalign{}
Serviço & Disponibilidade \\
\midrule\noalign{}
\endhead
\bottomrule\noalign{}
\endlastfoot
USF /\allowbreak  Posto de Saúde & sim \\
Hospital Regional & não \\
IPS (seguro social) & sim \\
Farmácia & sim \\
Distância ao hospital de referência & \textasciitilde47 km
(Concepción) \\
\end{longtable}
}

\textbf{Principais estabelecimentos:} USF local; referencia hospitalar
em Concepción

\textbf{Observação para imigrantes:} Atendimento primario local ou em
raio curto; casos de maior complexidade seguem para Concepción.
Cobertura privada continua recomendada para especialidades e urgências
de maior complexidade.
\newpage
\secaoDiagnostico{San Alfredo}{\mapaDistrito{0110}}{San Alfredo é uma excelente opção para quem busca baixíssima densidade
demográfica com acesso rodoviário moderno (PY22). O GSS de 6.4 reflete
um equilíbrio entre ótimos recursos naturais e os desafios de segurança
regional do departamento de Concepción.

\subsubsection{Por que sim}

\subsubsection{Por que não}

\subsubsection{Recomendação
Técnica}

Indicado para estabelecimentos rurais de médio/\allowbreak grande porte
(fazendas/\allowbreak silvicultura). Exige conformidade legal rigorosa para
estrangeiros devido à zona de fronteira. Referência atual de score: GSS
6.4.

\begin{itemize}
\tightlist
\item
  \begin{enumerate}
  \def\labelenumi{\arabic{enumi})}
  \tightlist
  \item
    Censo INE\footnote{\fineINE} 2022.
  \end{enumerate}
\item
  \begin{enumerate}
  \def\labelenumi{\arabic{enumi})}
  \setcounter{enumi}{1}
  \tightlist
  \item
    Relatórios de infraestrutura MOPC (Ruta PY22).
  \end{enumerate}
\item
  \begin{enumerate}
  \def\labelenumi{\arabic{enumi})}
  \setcounter{enumi}{2}
  \tightlist
  \item
    Mapa de sismicidade e riscos ambientais (SEN\footnote{\fineSEN}/\allowbreak DMH).
  \end{enumerate}
\end{itemize}}{dist:01_Concepción:San_Alfredo}
\begin{table}[ht!]
\small \begin{tabular}{p{3.0cm}p{0.8cm}p{6.0cm}} \toprule
\textbf{Critério} & \textbf{Nota} & \textbf{Justificativa} \\ \midrule
A - Ameaças Estratégicas & 6.0 & Nó logístico na Ruta PY22; zona de transição estratégica \\
B - Risco Social & 6.5 & Densidade baixíssima compensada pela presença de insurgência regional \\
C - Riscos Naturais & 6.5 & Risco de inundação moderado; geologia estável \\
D - Autossuficiência & 7.5 & Solo fértil para pastagem/\allowbreak silvicultura; boa disponibilidade hídrica \\
E - Institucional & 5.0 & Município recente; restrições legais de fronteira para brasileiros \\
\midrule \textbf{GSS FINAL} & \textbf{6.4} & \textbf{Moderadamente Seguro (Bom potencial para relocação com foco em pecuária ou silvicultura).} \\ \bottomrule \end{tabular}
\end{table}
\subsubsection{Vulnerabilidades e
Mitigação}

\begin{itemize}
\tightlist
\item
  \textbf{Isolamento Geográfico Total:} Risco de ficar sem suprimentos
  terrestres por meses (Mitigação: posse obrigatória de embarcação a
  motor robusta e, idealmente, aeronave leve; manutenção de estoque de
  segurança para 180 dias).
\item
  \textbf{Inexistência de Suporte Vital:} Distância crítica de hospitais
  cirúrgicos (Mitigação: manutenção de posto médico privado avançado e
  sistema de telemedicina satelital).
\item
  \textbf{Fragilidade Elétrica Extrema:} Quedas frequentes da rede
  (Mitigação: instalação de sistemas solares fotovoltaicos potentes
  off-grid com baterias de lítio).
\end{itemize}

\subsubsection{Dossiê de Campo}
\begin{description}[style=nextline,leftmargin=0.5cm,font=\bfseries]
\item[Coordenadas]  20°13'S 58°10'W.
\item[Topografia]  Relevo plano, zona de transição entre o Chaco Boreal e o Pantanal paraguaio. Altitude \textasciitilde 80m. Área massiva de \textasciitilde 36.513 km². Geologicamente estável, situada em uma depressão crustal que originou os humedales.
\item[Alvos Estratégicos]  Porto fluvial no Rio Paraguai; Base Naval de Bahía Negra (controle de fronteira fluvial); Estação Biológica Los Tres Gigantes (ativo ambiental); Ponto de fronteira tripartite (Paraguai, Bolívia e Brasil). Ausência total de infraestrutura industrial pesada, oferecendo o nível máximo de invisibilidade estratégica e proteção por isolamento geográfico absoluto.
\item[Fallout]  Ventos predominantes do Norte (NE) e Leste. Risco de fallout nulo de alvos nacionais; posição isolada no extremo norte do território.
\item[População]  2.768 habitantes (Censo 2022 Final). O distrito menos povoado de Alto Paraguay, composto por comunidades tradicionais, indígenas (Yshir) e militares.
\item[Densidade]  \textasciitilde 0,07 hab/\allowbreak km² (A menor densidade demográfica do Paraguai). Garante privacidade absoluta e ausência total de riscos sociais de massa urbana.
\item[Vias de Acesso]  Isolamento Crítico. Não há acesso asfáltico. Dependência total de estradas de terra precárias que se tornam intransitáveis por períodos de 40 a 100 dias durante chuvas intensas. A principal artéria vital é o Rio Paraguai (navegação fluvial). Existe um projeto de asfaltamento de 115 km até a Bioceânica, ainda sem financiamento pleno.
\item[Serviços]  Infraestrutura institucional e comercial mínima. Dependência de serviços de saúde da Marinha ou deslocamento aéreo/\allowbreak fluvial para centros maiores. A vila urbana é um micro núcleo de subsistência e suporte à fronteira.
\item[Custo de Vida]  Médio-alto devido ao custo do frete logístico; economia baseada na pecuária extensiva e ecoturismo.
\item[Preço da Terra]  US\$ 150-450/\allowbreak ha (campos brutos de interior); US\$ 500-950/\allowbreak ha (áreas com infraestrutura pecuária básica).
\item[Hidrologia]  Risco de Inundação Crítico. Distrito inserido no complexo do Pantanal. Vulnerabilidade extrema a cheias sazonais do Rio Paraguai e Rio Negro que isolam a localidade por meses. Solo hidromórfico (Solontz) com drenagem lentíssima.
\item[Clima]  Tropical de Savana com influência de humedales. Verões com calor extremo (>42°C) e alta umidade; invernos amenos com invasões de frentes frias austrais.
\item[Energia]  Rede nacional (Itaipu) via longas linhas rurais, altamente instável. Uso obrigatório de sistemas fotovoltaicos e geradores a diesel (suprimento fluvial).
\item[Água]  Abundância hídrica superficial (Rio Paraguai), exigindo tratamento. Água subterrânea frequentemente salobra ou profunda.
\item[Qualidade do Solo]  Solo arenoso e argiloso; excelente aptidão para pastagens naturais e pecuária de corte. Agricultura mecanizada inviável sem projetos massivos de engenharia hídrica.
\item[Recursos Locais]  Grande excedente de gado bovino e recursos pesqueiros. Elevadíssimo potencial para autossuficiência de proteína animal para a ínfima carga populacional.
\item[Segurança]  Zona de fronteira tripartite pacífica, mas sob vigilância militar permanente (Marinha). Criminalidade urbana comum inexistente. Presença marcante de comunidades indígenas originárias. Ideal para quem busca desaparecimento estratégico absoluto, protegida pela hostilidade da geografia.
\item[Leis Local: Município consolidado. Restrição de Fronteira (Lei 2532/\allowbreak 05)]  Aplicável a terras rurais na faixa de 50km da fronteira boliviana e brasileira.
\end{description}
\textbf{Fonte climática:} NASA POWER Climatology API (período 2001-2020)

\textbf{Fonte luminosa:} estimativa world atlas

\paragraph{Irradiação Solar
(kWh/\allowbreak m²/\allowbreak dia)}

{\def\LTcaptype{none} % do not increment counter
\begin{longtable}[]{@{}
  >{\raggedright\arraybackslash}p{(\linewidth - 24\tabcolsep) * \real{0.0746}}
  >{\raggedright\arraybackslash}p{(\linewidth - 24\tabcolsep) * \real{0.0746}}
  >{\raggedright\arraybackslash}p{(\linewidth - 24\tabcolsep) * \real{0.0746}}
  >{\raggedright\arraybackslash}p{(\linewidth - 24\tabcolsep) * \real{0.0746}}
  >{\raggedright\arraybackslash}p{(\linewidth - 24\tabcolsep) * \real{0.0746}}
  >{\raggedright\arraybackslash}p{(\linewidth - 24\tabcolsep) * \real{0.0746}}
  >{\raggedright\arraybackslash}p{(\linewidth - 24\tabcolsep) * \real{0.0746}}
  >{\raggedright\arraybackslash}p{(\linewidth - 24\tabcolsep) * \real{0.0746}}
  >{\raggedright\arraybackslash}p{(\linewidth - 24\tabcolsep) * \real{0.0746}}
  >{\raggedright\arraybackslash}p{(\linewidth - 24\tabcolsep) * \real{0.0746}}
  >{\raggedright\arraybackslash}p{(\linewidth - 24\tabcolsep) * \real{0.0746}}
  >{\raggedright\arraybackslash}p{(\linewidth - 24\tabcolsep) * \real{0.0746}}
  >{\raggedright\arraybackslash}p{(\linewidth - 24\tabcolsep) * \real{0.1045}}@{}}
\toprule\noalign{}
\begin{minipage}[b]{\linewidth}\raggedright
Jan
\end{minipage} & \begin{minipage}[b]{\linewidth}\raggedright
Fev
\end{minipage} & \begin{minipage}[b]{\linewidth}\raggedright
Mar
\end{minipage} & \begin{minipage}[b]{\linewidth}\raggedright
Abr
\end{minipage} & \begin{minipage}[b]{\linewidth}\raggedright
Mai
\end{minipage} & \begin{minipage}[b]{\linewidth}\raggedright
Jun
\end{minipage} & \begin{minipage}[b]{\linewidth}\raggedright
Jul
\end{minipage} & \begin{minipage}[b]{\linewidth}\raggedright
Ago
\end{minipage} & \begin{minipage}[b]{\linewidth}\raggedright
Set
\end{minipage} & \begin{minipage}[b]{\linewidth}\raggedright
Out
\end{minipage} & \begin{minipage}[b]{\linewidth}\raggedright
Nov
\end{minipage} & \begin{minipage}[b]{\linewidth}\raggedright
Dez
\end{minipage} & \begin{minipage}[b]{\linewidth}\raggedright
Média
\end{minipage} \\
\midrule\noalign{}
\endhead
\bottomrule\noalign{}
\endlastfoot
6.35 & 5.85 & 5.61 & 4.86 & 3.73 & 3.36 & 3.71 & 4.40 & 4.74 & 5.38 &
6.15 & 6.29 & \textbf{5.04} \\
\end{longtable}
}

\textbf{Inclinação solar recomendada:} 20° N (anual) · 30° N (inverno
jun-ago) · 10° N (verão nov-jan)

\paragraph{Precipitação (mm/\allowbreak mês)}

{\def\LTcaptype{none} % do not increment counter
\begin{longtable}[]{@{}
  >{\raggedright\arraybackslash}p{(\linewidth - 24\tabcolsep) * \real{0.0704}}
  >{\raggedright\arraybackslash}p{(\linewidth - 24\tabcolsep) * \real{0.0704}}
  >{\raggedright\arraybackslash}p{(\linewidth - 24\tabcolsep) * \real{0.0704}}
  >{\raggedright\arraybackslash}p{(\linewidth - 24\tabcolsep) * \real{0.0704}}
  >{\raggedright\arraybackslash}p{(\linewidth - 24\tabcolsep) * \real{0.0704}}
  >{\raggedright\arraybackslash}p{(\linewidth - 24\tabcolsep) * \real{0.0704}}
  >{\raggedright\arraybackslash}p{(\linewidth - 24\tabcolsep) * \real{0.0704}}
  >{\raggedright\arraybackslash}p{(\linewidth - 24\tabcolsep) * \real{0.0704}}
  >{\raggedright\arraybackslash}p{(\linewidth - 24\tabcolsep) * \real{0.0704}}
  >{\raggedright\arraybackslash}p{(\linewidth - 24\tabcolsep) * \real{0.0704}}
  >{\raggedright\arraybackslash}p{(\linewidth - 24\tabcolsep) * \real{0.0704}}
  >{\raggedright\arraybackslash}p{(\linewidth - 24\tabcolsep) * \real{0.0704}}
  >{\raggedright\arraybackslash}p{(\linewidth - 24\tabcolsep) * \real{0.1549}}@{}}
\toprule\noalign{}
\begin{minipage}[b]{\linewidth}\raggedright
Jan
\end{minipage} & \begin{minipage}[b]{\linewidth}\raggedright
Fev
\end{minipage} & \begin{minipage}[b]{\linewidth}\raggedright
Mar
\end{minipage} & \begin{minipage}[b]{\linewidth}\raggedright
Abr
\end{minipage} & \begin{minipage}[b]{\linewidth}\raggedright
Mai
\end{minipage} & \begin{minipage}[b]{\linewidth}\raggedright
Jun
\end{minipage} & \begin{minipage}[b]{\linewidth}\raggedright
Jul
\end{minipage} & \begin{minipage}[b]{\linewidth}\raggedright
Ago
\end{minipage} & \begin{minipage}[b]{\linewidth}\raggedright
Set
\end{minipage} & \begin{minipage}[b]{\linewidth}\raggedright
Out
\end{minipage} & \begin{minipage}[b]{\linewidth}\raggedright
Nov
\end{minipage} & \begin{minipage}[b]{\linewidth}\raggedright
Dez
\end{minipage} & \begin{minipage}[b]{\linewidth}\raggedright
Total/\allowbreak ano
\end{minipage} \\
\midrule\noalign{}
\endhead
\bottomrule\noalign{}
\endlastfoot
135 & 134 & 102 & 64 & 66 & 25 & 20 & 14 & 38 & 85 & 103 & 139 &
\textbf{930 mm} \\
\end{longtable}
}

\paragraph{Poluição Luminosa}

{\def\LTcaptype{none} % do not increment counter
\begin{longtable}[]{@{}ll@{}}
\toprule\noalign{}
Parâmetro & Valor \\
\midrule\noalign{}
\endhead
\bottomrule\noalign{}
\endlastfoot
Escala Bortle & 1 --- Céu verdadeiramente escuro \\
Radiância artificial & 0.1 nW/\allowbreak cm²/\allowbreak sr \\
\end{longtable}
}
\paragraph{Solo (SoilGrids 2.0, média ponderada 0--30
cm)}

{\def\LTcaptype{none} % do not increment counter
\begin{longtable}[]{@{}ll@{}}
\toprule\noalign{}
Parâmetro & Valor \\
\midrule\noalign{}
\endhead
\bottomrule\noalign{}
\endlastfoot
pH (H$_2$O) & 6.15 \\
Carbono orgânico (SOC) & 12.92 g/\allowbreak kg \\
Argila & 34.32 \% \\
Areia & 34.35 \% \\
Silte (calc.) & 31.3 \% \\
Densidade aparente & 1.33 g/\allowbreak cm³ \\
\textbf{Aptidão agrícola} & \textbf{Alta} \\
\end{longtable}
}

\subsubsection{Indicadores Sociais}

\textbf{Fontes:} PNUD 2020, INE\footnote{\fineINE} Censo 2022, INE\footnote{\fineINE} EPHC 2023, Ministerio
Público 2024\\
\textbf{Nota:} Valores marcados como \emph{(dept.)} referem-se ao
departamento de Alto Paraguay; valores \emph{(dist.)} são específicos
deste distrito. Dados marcados \emph{(est.)} são estimativas.

{\def\LTcaptype{none} % do not increment counter
\begin{longtable}[]{@{}
  >{\raggedright\arraybackslash}p{(\linewidth - 4\tabcolsep) * \real{0.4231}}
  >{\raggedright\arraybackslash}p{(\linewidth - 4\tabcolsep) * \real{0.2692}}
  >{\raggedright\arraybackslash}p{(\linewidth - 4\tabcolsep) * \real{0.3077}}@{}}
\toprule\noalign{}
\begin{minipage}[b]{\linewidth}\raggedright
Indicador
\end{minipage} & \begin{minipage}[b]{\linewidth}\raggedright
Valor
\end{minipage} & \begin{minipage}[b]{\linewidth}\raggedright
Âmbito
\end{minipage} \\
\midrule\noalign{}
\endhead
\bottomrule\noalign{}
\endlastfoot
IDH (2020) & 0,619 & dept. \\
Ranking IDH nacional & 18/\allowbreak 18 (ou 17/\allowbreak 18, disputado com Caazapá) &
dept. \\
Esperança de vida & 68,4 anos & dept. \\
Escolaridade média & 7.2 anos & dist. \\
RNB per capita (USD PPA) & 4.900 & dept. \\
Pobreza monetária (\%) & 38,7\% & dept. \\
Pobreza extrema (\%) & estimativa 14,2\% & dept. \\
Índice de Gini & estimativa 0,478 & dept. \\
Acesso a água potável (\%) & estimativa 34,5\% & dept. \\
Acesso a saneamento (\%) & estimativa 31,2\% & dept. \\
Taxa de homicídios (est., /\allowbreak 100k hab) & \textasciitilde2--4 (estimativa,
abaixo da média nacional) & dept. \\
Índice de segurança & alto & dept. \\
População (Censo 2022) & 2.768 habitantes & dist. \\
Idade mediana & N/\allowbreak D & dist. \\
Taxa de fecundidade & N/\allowbreak D & dist. \\
\end{longtable}
}
\subsubsection{Combustível}

\textbf{Referência:} PETROPAR /\allowbreak  postos locais (2024) \textbf{Tipo de
localidade:} Interior

{\def\LTcaptype{none} % do not increment counter
\begin{longtable}[]{@{}lll@{}}
\toprule\noalign{}
Combustível & USD/\allowbreak litro & Gs/\allowbreak litro (aprox.) \\
\midrule\noalign{}
\endhead
\bottomrule\noalign{}
\endlastfoot
Gasolina 93 oct & 1.18 & 8,732 \\
Gasolina 97 oct (premium) & 1.28 & 9,472 \\
Diesel & 1.10 & 8,140 \\
\end{longtable}
}

\begin{quote}
Preços podem variar ±5\% conforme posto e sazonalidade. Chaco e interior
remoto apresentam maior variação.
\end{quote}

\subsubsection{Cobertura Celular}

\textbf{Fonte:} CONATEL PY /\allowbreak  operadoras (2024)

{\def\LTcaptype{none} % do not increment counter
\begin{longtable}[]{@{}ll@{}}
\toprule\noalign{}
Parâmetro & Valor \\
\midrule\noalign{}
\endhead
\bottomrule\noalign{}
\endlastfoot
Cobertura 4G população (dept.) & 35\% \\
Cobertura 4G área rural & 15\% \\
Melhor operadora & Tigo \\
Qualidade rural & limitada \\
\end{longtable}
}

\begin{quote}
Para áreas rurais fora do núcleo urbano, recomenda-se chip Tigo como
principal e Personal como backup.
\end{quote}

\subsubsection{Internet}

\textbf{Fonte:} CONATEL /\allowbreak  Speedtest Ookla (2024)

{\def\LTcaptype{none} % do not increment counter
\begin{longtable}[]{@{}ll@{}}
\toprule\noalign{}
Parâmetro & Valor \\
\midrule\noalign{}
\endhead
\bottomrule\noalign{}
\endlastfoot
Velocidade média download & 15 Mbps \\
Domicílios com internet (dept.) & 22\% \\
Tecnologia predominante & satélite \\
Opção rural & Starlink disponível (\textasciitilde USD 44/\allowbreak mês) \\
\end{longtable}
}

\subsubsection{Mercado Imobiliário e Terra
Rural}

\textbf{Fonte:} INDERT /\allowbreak  Clasificados.com.py (2024)

{\def\LTcaptype{none} % do not increment counter
\begin{longtable}[]{@{}ll@{}}
\toprule\noalign{}
Tipo & Referência \\
\midrule\noalign{}
\endhead
\bottomrule\noalign{}
\endlastfoot
Terra agrícola alta prod. (USD/\allowbreak ha) & 600 \\
Imóvel urbano (USD/\allowbreak m²) & 400 \\
Aluguel 2 quartos (USD/\allowbreak mês) & 160 \\
\end{longtable}
}

\begin{quote}
Valores de referência departamental. Localidades menores podem ter
preços 20--40\% abaixo da capital departamental.
\end{quote}

\subsubsection{Saúde}

\textbf{Fonte:} MSPBS /\allowbreak  IPS Paraguay (2024-2026), consolidação
departamental e proxy local

{\def\LTcaptype{none} % do not increment counter
\begin{longtable}[]{@{}ll@{}}
\toprule\noalign{}
Serviço & Disponibilidade \\
\midrule\noalign{}
\endhead
\bottomrule\noalign{}
\endlastfoot
USF /\allowbreak  Posto de Saúde & sim \\
Hospital Regional & não \\
IPS (seguro social) & não \\
Farmácia & sim \\
Distância ao hospital de referência & \textasciitilde119 km (Fuerte
Olimpo) \\
\end{longtable}
}

\textbf{Principais estabelecimentos:} USF local; referencia hospitalar
em Fuerte Olimpo

\textbf{Observação para imigrantes:} Atendimento primario local ou em
raio curto; casos de maior complexidade seguem para Fuerte Olimpo.
Cobertura privada continua recomendada para especialidades e urgências
de maior complexidade.
\newpage
\secaoDiagnostico{San Carlos del Apa}{\mapaDistrito{0105}}{Capitán Carmelo Peralta (Alto Paraguay) é um hub de fronteira em
transição radical. Deixou de ser uma zona isolada para se tornar o
epicentro da Rota Bioceânica, o que traz oportunidades massivas de
autossuficiência proteica e valorização, mas também riscos estratégicos
elevados por ser um nó logístico internacional.

\subsubsection{Por que sim}

\subsubsection{Por que não}

\subsubsection{Pre-condicoes Minimas de
Decisao}

\subsubsection{Recomendação
Técnica}

Uso recomendado para perfis que buscam valorização estratégica e
autossuficiência em áreas de baixa densidade, cientes do papel
geopolítico do hub. Referência atual de score: GSS 6.8; risco
predominante: moderado (hidrológico e estratégico).

\begin{itemize}
\tightlist
\item
  \begin{enumerate}
  \def\labelenumi{\arabic{enumi})}
  \tightlist
  \item
    Delimitacao territorial validada por georreferencia institucional.
  \end{enumerate}
\item
  \begin{enumerate}
  \def\labelenumi{\arabic{enumi})}
  \setcounter{enumi}{1}
  \tightlist
  \item
    População municipal referenciada por Censo 2022.
  \end{enumerate}
\item
  \begin{enumerate}
  \def\labelenumi{\arabic{enumi})}
  \setcounter{enumi}{2}
  \tightlist
  \item
    Indicadores de infraestrutura e rodovias PY15.
  \end{enumerate}
\item
  \begin{enumerate}
  \def\labelenumi{\arabic{enumi})}
  \setcounter{enumi}{3}
  \tightlist
  \item
    Riscos climáticos e hidrológicos.
  \end{enumerate}
\item
  \begin{enumerate}
  \def\labelenumi{\arabic{enumi})}
  \setcounter{enumi}{4}
  \tightlist
  \item
    Informações de custo de energia e matriz.
  \end{enumerate}
\item
  \begin{enumerate}
  \def\labelenumi{\arabic{enumi})}
  \setcounter{enumi}{5}
  \tightlist
  \item
    Restrições legais de fronteira - Lei 2532/\allowbreak 05.
  \end{enumerate}
\end{itemize}}{dist:01_Concepción:San_Carlos_del_Apa}
\begin{table}[ht!]
\small \begin{tabular}{p{3.0cm}p{0.8cm}p{6.0cm}} \toprule
\textbf{Critério} & \textbf{Nota} & \textbf{Justificativa} \\ \midrule
A - Ameaças Estratégicas & 6.0 & Ponto estratégico de fronteira, mas remoto o suficiente para não ser alvo primário global \\
B - Risco Social & 8.5 & Densidade demográfica entre as menores do país; risco social baixíssimo \\
C - Riscos Naturais & 6.5 & Geologia estável no Chaco; vulnerabilidade hídrica sazonal \\
D - Autossuficiência & 8.0 & Elevado potencial de autossuficiência proteica e hub de suprimentos \\
E - Institucional & 7.4 & Forte suporte governamental; restrição legal de fronteira para estrangeiros \\
\midrule \textbf{GSS FINAL} & \textbf{6.8} & \textbf{Moderadamente Seguro (Exige preparação e infraestrutura dedicada).} \\ \bottomrule \end{tabular}
\end{table}
\subsubsection{Vulnerabilidades e
Mitigação}

\begin{itemize}
\tightlist
\item
  \textbf{Exposicao Hidrologica:} Risco de inundação sazonal pelo Rio
  Paraguai (reforço de defesas costeiras e construção elevada
  recomendada).
\item
  \textbf{Isolamento Logístico:} Dependência de poucas rotas de saída; o
  Corredor Bioceânico é vital (monitorar tráfego e possíveis bloqueios).
\item
  \textbf{Segurança de Fronteira:} Proximidade com o Brasil exige
  vigilância contra fluxos migratórios descontrolados ou crises
  transfronteiriças.
\item
  \textbf{Autossuficiência Energética:} Baixa diversidade de fontes
  locais (potencializar microgeração solar para reduzir dependência da
  rede).
\end{itemize}

\subsubsection{Dossiê de Campo}
\begin{description}[style=nextline,leftmargin=0.5cm,font=\bfseries]
\item[Coordenadas]  21°41'S 57°54'W.
\item[Topografia]  Relevo plano (geologicamente estável, baixo risco sísmico), margeado pelo Rio Paraguai (fronteira com Brasil). Área de \textasciitilde 4.798 km².
\item[Alvos Estratégicos]  Ponte Bioceânica (em fase final de construção - conexão vital Paraguai-Brasil); Eixo da Ruta PY15 (Corredor Bioceânico); Porto fluvial; Proximidade com o Pantanal paraguaio. Localização estratégica absoluta como hub logístico internacional de carga (celulose, grãos e proteína).
\item[Fallout]  Ventos predominantes do Norte/\allowbreak Nordeste.
\item[População]  4.402 (Censo 2022 Final). População em rápido crescimento devido ao impacto das obras de infraestrutura. Presença de comunidades indígenas Ayoreo.
\item[Densidade]  \textasciitilde 0,92 hab/\allowbreak km² (Extremamente baixa, oferecendo isolamento geográfico massivo fora do núcleo urbano).
\item[Vias de Acesso]  Transformação radical com o Corredor Bioceânico pavimentado (PY15) e a nova ponte internacional para Porto Murtinho (Brasil). Risco elevado de saturação e controle militar em crises devido ao valor logístico da rota.
\item[Serviços]  Infraestrutura urbana e de serviços em fase de expansão acelerada. Dependência residual de centros maiores para saúde de alta complexidade.
\item[Custo de Vida]  Em ascensão devido à valorização imobiliária; energia elétrica barata (ANDE\footnote{\fineANDE} \textasciitilde US\$ 0,05/\allowbreak kWh).
\item[Preço da Terra]  US\$ 300-600/\allowbreak ha (bruta) a US\$ 1.200/\allowbreak ha (desenvolvida). Forte especulação em áreas logísticas.
\item[Hidrologia]  Risco de Inundação Sazonal. Localidade margeada pelo Rio Paraguai, vulnerável a cheias extraordinárias. Obras de defesa costeira em andamento. Risco de incêndios florestais em períodos de seca no Chaco.
\item[Clima]  Tropical Úmido/\allowbreak Seco. Verões com calor extremo. Déficit hídrico é o principal desafio agrícola.
\item[Sismicidade]  Risco baixo; região estável com eventos intraplaca raros e de baixa magnitude (<4.0 Richter).
\item[Energia]  Rede nacional (Itaipu) em fase de reforço para atender ao novo hub logístico.
\item[Água]  Rio Paraguai e sistemas de tratamento local.
\item[Qualidade do Solo]  Solos aluviais ricos em fósforo; desafios de drenagem e salinidade. Potencial agrícola em expansão (soja, milho, algodão) com manejo adequado.
\item[Recursos Locais]  Pecuária de corte intensiva e pesca. Elevadíssimo potencial para autossuficiência a nível de proteína animal devido à baixa densidade e vastidão territorial. Localidade estratégica para suprimento de longo prazo se a logística bioceânica estiver operacional.
\item[Segurança]  Zona de fronteira em transformação. Baixa criminalidade urbana violenta tradicional, mas com desafios crescentes ligados ao fluxo internacional de carga. Coesão social desafiada pelo rápido crescimento populacional e migratório.
\item[Leis Local: Município consolidado e foco prioritário de investimentos nacionais. Restrição de Fronteira (Lei 2532/\allowbreak 05)]  Estrangeiros de países limítrofes não podem possuir terras em faixa de 50km da fronteira sem decreto de interesse público.
\end{description}
Coordenadas: -21.04°S /\allowbreak  -57.87°W. Acesso em 2026-03-20.

\textbf{Irradiância solar global (kWh/\allowbreak m²/\allowbreak dia):}

{\def\LTcaptype{none} % do not increment counter
\begin{longtable}[]{@{}
  >{\raggedright\arraybackslash}p{(\linewidth - 22\tabcolsep) * \real{0.0833}}
  >{\raggedright\arraybackslash}p{(\linewidth - 22\tabcolsep) * \real{0.0833}}
  >{\raggedright\arraybackslash}p{(\linewidth - 22\tabcolsep) * \real{0.0833}}
  >{\raggedright\arraybackslash}p{(\linewidth - 22\tabcolsep) * \real{0.0833}}
  >{\raggedright\arraybackslash}p{(\linewidth - 22\tabcolsep) * \real{0.0833}}
  >{\raggedright\arraybackslash}p{(\linewidth - 22\tabcolsep) * \real{0.0833}}
  >{\raggedright\arraybackslash}p{(\linewidth - 22\tabcolsep) * \real{0.0833}}
  >{\raggedright\arraybackslash}p{(\linewidth - 22\tabcolsep) * \real{0.0833}}
  >{\raggedright\arraybackslash}p{(\linewidth - 22\tabcolsep) * \real{0.0833}}
  >{\raggedright\arraybackslash}p{(\linewidth - 22\tabcolsep) * \real{0.0833}}
  >{\raggedright\arraybackslash}p{(\linewidth - 22\tabcolsep) * \real{0.0833}}
  >{\raggedright\arraybackslash}p{(\linewidth - 22\tabcolsep) * \real{0.0833}}@{}}
\toprule\noalign{}
\begin{minipage}[b]{\linewidth}\raggedright
Jan
\end{minipage} & \begin{minipage}[b]{\linewidth}\raggedright
Fev
\end{minipage} & \begin{minipage}[b]{\linewidth}\raggedright
Mar
\end{minipage} & \begin{minipage}[b]{\linewidth}\raggedright
Abr
\end{minipage} & \begin{minipage}[b]{\linewidth}\raggedright
Mai
\end{minipage} & \begin{minipage}[b]{\linewidth}\raggedright
Jun
\end{minipage} & \begin{minipage}[b]{\linewidth}\raggedright
Jul
\end{minipage} & \begin{minipage}[b]{\linewidth}\raggedright
Ago
\end{minipage} & \begin{minipage}[b]{\linewidth}\raggedright
Set
\end{minipage} & \begin{minipage}[b]{\linewidth}\raggedright
Out
\end{minipage} & \begin{minipage}[b]{\linewidth}\raggedright
Nov
\end{minipage} & \begin{minipage}[b]{\linewidth}\raggedright
Dez
\end{minipage} \\
\midrule\noalign{}
\endhead
\bottomrule\noalign{}
\endlastfoot
6.43 & 5.94 & 5.66 & 4.87 & 3.73 & 3.33 & 3.58 & 4.35 & 4.83 & 5.51 &
6.31 & 6.42 \\
\end{longtable}
}

Média anual: 5.08 kWh/\allowbreak m²/\allowbreak dia. Mínimo: Jun (3.33) \textbar{} Máximo: Jan
(6.43).

\textbf{Precipitação (mm/\allowbreak mês → mm/\allowbreak ano \textasciitilde1.005):}

{\def\LTcaptype{none} % do not increment counter
\begin{longtable}[]{@{}
  >{\raggedright\arraybackslash}p{(\linewidth - 22\tabcolsep) * \real{0.0833}}
  >{\raggedright\arraybackslash}p{(\linewidth - 22\tabcolsep) * \real{0.0833}}
  >{\raggedright\arraybackslash}p{(\linewidth - 22\tabcolsep) * \real{0.0833}}
  >{\raggedright\arraybackslash}p{(\linewidth - 22\tabcolsep) * \real{0.0833}}
  >{\raggedright\arraybackslash}p{(\linewidth - 22\tabcolsep) * \real{0.0833}}
  >{\raggedright\arraybackslash}p{(\linewidth - 22\tabcolsep) * \real{0.0833}}
  >{\raggedright\arraybackslash}p{(\linewidth - 22\tabcolsep) * \real{0.0833}}
  >{\raggedright\arraybackslash}p{(\linewidth - 22\tabcolsep) * \real{0.0833}}
  >{\raggedright\arraybackslash}p{(\linewidth - 22\tabcolsep) * \real{0.0833}}
  >{\raggedright\arraybackslash}p{(\linewidth - 22\tabcolsep) * \real{0.0833}}
  >{\raggedright\arraybackslash}p{(\linewidth - 22\tabcolsep) * \real{0.0833}}
  >{\raggedright\arraybackslash}p{(\linewidth - 22\tabcolsep) * \real{0.0833}}@{}}
\toprule\noalign{}
\begin{minipage}[b]{\linewidth}\raggedright
Jan
\end{minipage} & \begin{minipage}[b]{\linewidth}\raggedright
Fev
\end{minipage} & \begin{minipage}[b]{\linewidth}\raggedright
Mar
\end{minipage} & \begin{minipage}[b]{\linewidth}\raggedright
Abr
\end{minipage} & \begin{minipage}[b]{\linewidth}\raggedright
Mai
\end{minipage} & \begin{minipage}[b]{\linewidth}\raggedright
Jun
\end{minipage} & \begin{minipage}[b]{\linewidth}\raggedright
Jul
\end{minipage} & \begin{minipage}[b]{\linewidth}\raggedright
Ago
\end{minipage} & \begin{minipage}[b]{\linewidth}\raggedright
Set
\end{minipage} & \begin{minipage}[b]{\linewidth}\raggedright
Out
\end{minipage} & \begin{minipage}[b]{\linewidth}\raggedright
Nov
\end{minipage} & \begin{minipage}[b]{\linewidth}\raggedright
Dez
\end{minipage} \\
\midrule\noalign{}
\endhead
\bottomrule\noalign{}
\endlastfoot
122 & 134 & 100 & 78 & 80 & 40 & 26 & 16 & 41 & 99 & 141 & 128 \\
\end{longtable}
}

Estação chuvosa Nov-Fev; seca Jun-Set (mínimo absoluto em Ago). Regime
sub-úmido do Chaco Boreal.

\textbf{Inclinação solar recomendada:} 21° Norte (anual); 31° no inverno
(Jun-Ago); 11° no verão (Nov-Jan). Base: latitude -21.04°S.
\paragraph{Solo (SoilGrids 2.0, média ponderada 0--30
cm)}

{\def\LTcaptype{none} % do not increment counter
\begin{longtable}[]{@{}ll@{}}
\toprule\noalign{}
Parâmetro & Valor \\
\midrule\noalign{}
\endhead
\bottomrule\noalign{}
\endlastfoot
pH (H$_2$O) & 6.25 \\
Carbono orgânico (SOC) & 13.4 g/\allowbreak kg \\
Argila & 31.66 \% \\
Areia & 36.19 \% \\
Silte (calc.) & 32.2 \% \\
Densidade aparente & 1.33 g/\allowbreak cm³ \\
\textbf{Aptidão agrícola} & \textbf{Alta} \\
\end{longtable}
}

\subsubsection{Indicadores Sociais}

\textbf{Fontes:} PNUD 2020, INE\footnote{\fineINE} Censo 2022, INE\footnote{\fineINE} EPHC 2023, Ministerio
Público 2024\\
\textbf{Nota:} Valores marcados como \emph{(dept.)} referem-se ao
departamento de Alto Paraguay; valores \emph{(dist.)} são específicos
deste distrito. Dados marcados \emph{(est.)} são estimativas.

{\def\LTcaptype{none} % do not increment counter
\begin{longtable}[]{@{}
  >{\raggedright\arraybackslash}p{(\linewidth - 4\tabcolsep) * \real{0.4231}}
  >{\raggedright\arraybackslash}p{(\linewidth - 4\tabcolsep) * \real{0.2692}}
  >{\raggedright\arraybackslash}p{(\linewidth - 4\tabcolsep) * \real{0.3077}}@{}}
\toprule\noalign{}
\begin{minipage}[b]{\linewidth}\raggedright
Indicador
\end{minipage} & \begin{minipage}[b]{\linewidth}\raggedright
Valor
\end{minipage} & \begin{minipage}[b]{\linewidth}\raggedright
Âmbito
\end{minipage} \\
\midrule\noalign{}
\endhead
\bottomrule\noalign{}
\endlastfoot
IDH (2020) & 0,619 & dept. \\
Ranking IDH nacional & 18/\allowbreak 18 (ou 17/\allowbreak 18, disputado com Caazapá) &
dept. \\
Esperança de vida & 68,4 anos & dept. \\
Escolaridade média & 8.1 anos & dist. \\
RNB per capita (USD PPA) & 4.900 & dept. \\
Pobreza monetária (\%) & 38,7\% & dept. \\
Pobreza extrema (\%) & estimativa 14,2\% & dept. \\
Índice de Gini & estimativa 0,478 & dept. \\
Acesso a água potável (\%) & estimativa 34,5\% & dept. \\
Acesso a saneamento (\%) & estimativa 31,2\% & dept. \\
Taxa de homicídios (est., /\allowbreak 100k hab) & \textasciitilde2--4 (estimativa,
abaixo da média nacional) & dept. \\
Índice de segurança & alto & dept. \\
População (Censo 2022) & 3.471 habitantes & dist. \\
Idade mediana & N/\allowbreak D & dist. \\
Taxa de fecundidade & N/\allowbreak D & dist. \\
\end{longtable}
}
\subsubsection{Combustível}

\textbf{Referência:} PETROPAR /\allowbreak  postos locais (2024) \textbf{Tipo de
localidade:} Capital departamental

{\def\LTcaptype{none} % do not increment counter
\begin{longtable}[]{@{}lll@{}}
\toprule\noalign{}
Combustível & USD/\allowbreak litro & Gs/\allowbreak litro (aprox.) \\
\midrule\noalign{}
\endhead
\bottomrule\noalign{}
\endlastfoot
Gasolina 93 oct & 1.18 & 8,732 \\
Gasolina 97 oct (premium) & 1.28 & 9,472 \\
Diesel & 1.10 & 8,140 \\
\end{longtable}
}

\begin{quote}
Preços podem variar ±5\% conforme posto e sazonalidade. Chaco e interior
remoto apresentam maior variação.
\end{quote}

\subsubsection{Cobertura Celular}

\textbf{Fonte:} CONATEL PY /\allowbreak  operadoras (2024)

{\def\LTcaptype{none} % do not increment counter
\begin{longtable}[]{@{}ll@{}}
\toprule\noalign{}
Parâmetro & Valor \\
\midrule\noalign{}
\endhead
\bottomrule\noalign{}
\endlastfoot
Cobertura 4G população (dept.) & 35\% \\
Cobertura 4G área rural & 15\% \\
Melhor operadora & Tigo \\
Qualidade rural & limitada \\
\end{longtable}
}

\begin{quote}
Para áreas rurais fora do núcleo urbano, recomenda-se chip Tigo como
principal e Personal como backup.
\end{quote}

\subsubsection{Internet}

\textbf{Fonte:} CONATEL /\allowbreak  Speedtest Ookla (2024)

{\def\LTcaptype{none} % do not increment counter
\begin{longtable}[]{@{}ll@{}}
\toprule\noalign{}
Parâmetro & Valor \\
\midrule\noalign{}
\endhead
\bottomrule\noalign{}
\endlastfoot
Velocidade média download & 15 Mbps \\
Domicílios com internet (dept.) & 22\% \\
Tecnologia predominante & satélite \\
Opção rural & Starlink disponível (\textasciitilde USD 44/\allowbreak mês) \\
\end{longtable}
}

\subsubsection{Mercado Imobiliário e Terra
Rural}

\textbf{Fonte:} INDERT /\allowbreak  Clasificados.com.py (2024)

{\def\LTcaptype{none} % do not increment counter
\begin{longtable}[]{@{}ll@{}}
\toprule\noalign{}
Tipo & Referência \\
\midrule\noalign{}
\endhead
\bottomrule\noalign{}
\endlastfoot
Terra agrícola alta prod. (USD/\allowbreak ha) & 600 \\
Imóvel urbano (USD/\allowbreak m²) & 459 \\
Aluguel 2 quartos (USD/\allowbreak mês) & 192 \\
\end{longtable}
}

\begin{quote}
Valores de referência departamental. Localidades menores podem ter
preços 20--40\% abaixo da capital departamental.
\end{quote}

\subsubsection{Saúde}

\textbf{Fonte:} MSPBS /\allowbreak  IPS Paraguay (2024-2026), consolidação
departamental e proxy local

{\def\LTcaptype{none} % do not increment counter
\begin{longtable}[]{@{}ll@{}}
\toprule\noalign{}
Serviço & Disponibilidade \\
\midrule\noalign{}
\endhead
\bottomrule\noalign{}
\endlastfoot
USF /\allowbreak  Posto de Saúde & sim \\
Hospital Regional & sim \\
IPS (seguro social) & sim \\
Farmácia & sim \\
Distância ao hospital de referência & local \\
\end{longtable}
}

\textbf{Principais estabelecimentos:} Hospital distrital /\allowbreak  regional de
Fuerte Olimpo; rede de USF e farmacias locais

\textbf{Observação para imigrantes:} Centro de referencia do
departamento. Acesso a atendimento primario e, em geral, a maior oferta
publica e privada da area. Cobertura privada continua recomendada para
especialidades e urgências de maior complexidade.
\newpage
\secaoDiagnostico{San Lazaro}{\mapaDistrito{0106}}{San Lázaro (Vallemí) é o coração industrial do Norte. Oferece uma
combinação única de recursos minerais estratégicos e conectividade
fluvial moderna, mas sua pontuação GSS é moderada devido ao seu papel
como alvo estratégico nacional e riscos hidrológicos.

\subsubsection{Por que sim}

\subsubsection{Por que não}

\subsubsection{Recomendação
Técnica}

Indicado para perfis com interesse em logística fluvial, mineração ou
apoio industrial. Para fins de refúgio isolado, os distritos vizinhos
(Itacuá ou San Carlos) oferecem menor densidade, mas com menor
infraestrutura. Referência atual de score: GSS 5.8.

\begin{itemize}
\tightlist
\item
  \begin{enumerate}
  \def\labelenumi{\arabic{enumi})}
  \tightlist
  \item
    Dados censitários INE\footnote{\fineINE} 2022.
  \end{enumerate}
\item
  \begin{enumerate}
  \def\labelenumi{\arabic{enumi})}
  \setcounter{enumi}{1}
  \tightlist
  \item
    Relatórios de infraestrutura e pavimentação MOPC.
  \end{enumerate}
\item
  \begin{enumerate}
  \def\labelenumi{\arabic{enumi})}
  \setcounter{enumi}{2}
  \tightlist
  \item
    Histórico de cheias e avisos SEN\footnote{\fineSEN}/\allowbreak DMH.
  \end{enumerate}
\end{itemize}}{dist:01_Concepción:San_Lazaro}
\begin{table}[ht!]
\small \begin{tabular}{p{3.0cm}p{0.8cm}p{6.0cm}} \toprule
\textbf{Critério} & \textbf{Nota} & \textbf{Justificativa} \\ \midrule
A - Ameaças Estratégicas & 5.0 & Hub industrial e logístico estratégico; alvo de interesse nacional \\
B - Risco Social & 6.5 & Densidade urbana moderada com controle estatal; localização de fronteira \\
C - Riscos Naturais & 6.0 & Geologia estável; risco de inundação fluvial cíclica \\
D - Autossuficiência & 6.5 & Recursos minerais abundantes e água; solo agrícola requer manejo \\
E - Institucional & 5.5 & Importância estratégica garante foco institucional, mas leis de fronteira são restritivas \\
\midrule \textbf{GSS FINAL} & \textbf{5.8} & \textbf{Moderadamente Seguro (Indicado para quem busca conexão com a indústria mineral e logística fluvial).} \\ \bottomrule \end{tabular}
\end{table}
\subsubsection{Vulnerabilidades e
Mitigação}

\begin{itemize}
\tightlist
\item
  \textbf{Alvo Industrial Estratégico:} A planta da INC é um ponto
  crítico em cenários de sabotagem ou crise nacional (Mitigação:
  residência fora do perímetro imediato da fábrica).
\item
  \textbf{Risco Hídrico:} Inundações sazonais podem afetar a logística
  fluvial e áreas baixas (Mitigação: construção em cotas elevadas,
  monitoramento de avisos do DMH).
\item
  \textbf{Restrições Legais:} Zona de fronteira exige conformidade
  rigorosa para investidores estrangeiros.
\item
  \textbf{Dependência de Serviços:} Necessidade de deslocamento para
  Concepción para serviços especializados.
\end{itemize}

\subsubsection{Dossiê de Campo}
\begin{description}[style=nextline,leftmargin=0.5cm,font=\bfseries]
\item[Coordenadas]  22°09'S 57°56'W.
\item[Topografia]  Localizada na confluência dos rios Paraguai e Apa. Terreno com elevações calcárias e complexos de cavernas (ex: Santa Caverna). Geologicamente estável.
\item[Alvos Estratégicos]  Industria Nacional del Cemento (INC) - planta vital e estratégica para a construção civil e infraestrutura nacional; Aeródromo de Vallemí (pista asfaltada); Porto de Vallemí (logística fluvial de carga pesada).
\item[Fallout]  Ventos predominantes do Norte/\allowbreak Nordeste.
\item[População]  11.192 (Censo 2022 Final).
\item[Densidade]  \textasciitilde 11,13 hab/\allowbreak km² (Moderada para a região, concentrada no núcleo urbano de Vallemí).
\item[Vias de Acesso]  Ruta PY22 (pavimentação moderna ligando a Concepción). Conexão fluvial estratégica pelos rios Paraguai e Apa. Balsa para o Brasil (Porto Murtinho) nas proximidades.
\item[Serviços]  Hospital da INC e centros de saúde básicos. Dependência de Concepción para alta complexidade.
\item[Custo de Vida]  Médio-baixo; influenciado pela atividade industrial e proximidade com o Brasil.
\item[Preço da Terra]  Lotes urbanos em Vallemí entre US\$ 15.000 e US\$ 40.000; áreas industriais/\allowbreak minerais com valorização elevada.
\item[Hidrologia]  Risco de Inundação Fluvial. Áreas baixas próximas à confluência dos rios Paraguai e Apa são vulneráveis a cheias cíclicas (eventos significativos em 2019 e 2023).
\item[Clima]  Tropical Úmido/\allowbreak Quente. Verões com temperaturas >40°C e tempestades severas.
\item[Energia]  Rede nacional (Itaipu). Subestação dedicada para a INC, conferindo certa prioridade de manutenção, mas sujeita a instabilidades regionais.
\item[Água]  Abundante (Rio Paraguai e Rio Apa).
\item[Qualidade do Solo]  Predomínio de formações calcárias; solo rico em minerais para construção e correção de solos, mas exige manejo para agricultura diversificada.
\item[Recursos Locais]  Maior polo produtor de cimento, cal e mármore do país. Pecuária e pesca comercial.
\item[Segurança]  Cidade industrial com forte presença estatal e corporativa. Zona de fronteira com o Brasil. Presença constante da FTC (Força de Tarefa Conjunta) na região norte do departamento. Taxa de homicídios regional de 14-19/\allowbreak 100k.
\item[Leis Local: Polo de mineração de importância nacional. Restrição de Fronteira (Lei 2532/\allowbreak 05)]  Proibição de compra de terras rurais por estrangeiros de países limítrofes (faixa de 50km).
\end{description}
\subsubsection{Indicadores Sociais}

\textbf{Fontes:} PNUD 2020, INE\footnote{\fineINE} Censo 2022, INE\footnote{\fineINE} EPHC 2023, Ministerio
Público 2024\\
\textbf{Nota:} Valores marcados como \emph{(dept.)} referem-se ao
departamento de Alto Paraguay; valores \emph{(dist.)} são específicos
deste distrito. Dados marcados \emph{(est.)} são estimativas.

{\def\LTcaptype{none} % do not increment counter
\begin{longtable}[]{@{}
  >{\raggedright\arraybackslash}p{(\linewidth - 4\tabcolsep) * \real{0.4231}}
  >{\raggedright\arraybackslash}p{(\linewidth - 4\tabcolsep) * \real{0.2692}}
  >{\raggedright\arraybackslash}p{(\linewidth - 4\tabcolsep) * \real{0.3077}}@{}}
\toprule\noalign{}
\begin{minipage}[b]{\linewidth}\raggedright
Indicador
\end{minipage} & \begin{minipage}[b]{\linewidth}\raggedright
Valor
\end{minipage} & \begin{minipage}[b]{\linewidth}\raggedright
Âmbito
\end{minipage} \\
\midrule\noalign{}
\endhead
\bottomrule\noalign{}
\endlastfoot
IDH (2020) & 0,619 & dept. \\
Ranking IDH nacional & 18/\allowbreak 18 (ou 17/\allowbreak 18, disputado com Caazapá) &
dept. \\
Esperança de vida & 68,4 anos & dept. \\
Escolaridade média & 7.7 anos & dist. \\
RNB per capita (USD PPA) & 4.900 & dept. \\
Pobreza monetária (\%) & 38,7\% & dept. \\
Pobreza extrema (\%) & estimativa 14,2\% & dept. \\
Índice de Gini & estimativa 0,478 & dept. \\
Acesso a água potável (\%) & estimativa 34,5\% & dept. \\
Acesso a saneamento (\%) & estimativa 31,2\% & dept. \\
Taxa de homicídios (est., /\allowbreak 100k hab) & \textasciitilde2--4 (estimativa,
abaixo da média nacional) & dept. \\
Índice de segurança & alto & dept. \\
População (Censo 2022) & 6.003 habitantes & dist. \\
Idade mediana & N/\allowbreak D & dist. \\
Taxa de fecundidade & N/\allowbreak D & dist. \\
\end{longtable}
}
\subsubsection{Combustível}

\textbf{Referência:} PETROPAR /\allowbreak  postos locais (2024) \textbf{Tipo de
localidade:} Interior

{\def\LTcaptype{none} % do not increment counter
\begin{longtable}[]{@{}lll@{}}
\toprule\noalign{}
Combustível & USD/\allowbreak litro & Gs/\allowbreak litro (aprox.) \\
\midrule\noalign{}
\endhead
\bottomrule\noalign{}
\endlastfoot
Gasolina 93 oct & 1.18 & 8,732 \\
Gasolina 97 oct (premium) & 1.28 & 9,472 \\
Diesel & 1.10 & 8,140 \\
\end{longtable}
}

\begin{quote}
Preços podem variar ±5\% conforme posto e sazonalidade. Chaco e interior
remoto apresentam maior variação.
\end{quote}

\subsubsection{Cobertura Celular}

\textbf{Fonte:} CONATEL PY /\allowbreak  operadoras (2024)

{\def\LTcaptype{none} % do not increment counter
\begin{longtable}[]{@{}ll@{}}
\toprule\noalign{}
Parâmetro & Valor \\
\midrule\noalign{}
\endhead
\bottomrule\noalign{}
\endlastfoot
Cobertura 4G população (dept.) & 35\% \\
Cobertura 4G área rural & 15\% \\
Melhor operadora & Tigo \\
Qualidade rural & limitada \\
\end{longtable}
}

\begin{quote}
Para áreas rurais fora do núcleo urbano, recomenda-se chip Tigo como
principal e Personal como backup.
\end{quote}

\subsubsection{Internet}

\textbf{Fonte:} CONATEL /\allowbreak  Speedtest Ookla (2024)

{\def\LTcaptype{none} % do not increment counter
\begin{longtable}[]{@{}ll@{}}
\toprule\noalign{}
Parâmetro & Valor \\
\midrule\noalign{}
\endhead
\bottomrule\noalign{}
\endlastfoot
Velocidade média download & 15 Mbps \\
Domicílios com internet (dept.) & 22\% \\
Tecnologia predominante & satélite \\
Opção rural & Starlink disponível (\textasciitilde USD 44/\allowbreak mês) \\
\end{longtable}
}

\subsubsection{Mercado Imobiliário e Terra
Rural}

\textbf{Fonte:} INDERT /\allowbreak  Clasificados.com.py (2024)

{\def\LTcaptype{none} % do not increment counter
\begin{longtable}[]{@{}ll@{}}
\toprule\noalign{}
Tipo & Referência \\
\midrule\noalign{}
\endhead
\bottomrule\noalign{}
\endlastfoot
Terra agrícola alta prod. (USD/\allowbreak ha) & 600 \\
Imóvel urbano (USD/\allowbreak m²) & 400 \\
Aluguel 2 quartos (USD/\allowbreak mês) & 160 \\
\end{longtable}
}

\begin{quote}
Valores de referência departamental. Localidades menores podem ter
preços 20--40\% abaixo da capital departamental.
\end{quote}

\subsubsection{Saúde}

\textbf{Fonte:} MSPBS /\allowbreak  IPS Paraguay (2024-2026), consolidação
departamental e proxy local

{\def\LTcaptype{none} % do not increment counter
\begin{longtable}[]{@{}ll@{}}
\toprule\noalign{}
Serviço & Disponibilidade \\
\midrule\noalign{}
\endhead
\bottomrule\noalign{}
\endlastfoot
USF /\allowbreak  Posto de Saúde & sim \\
Hospital Regional & sim \\
IPS (seguro social) & sim \\
Farmácia & sim \\
Distância ao hospital de referência & \textasciitilde174 km (Fuerte
Olimpo) \\
\end{longtable}
}

\textbf{Principais estabelecimentos:} USF local; referencia hospitalar
em Fuerte Olimpo

\textbf{Observação para imigrantes:} Atendimento primario local ou em
raio curto; casos de maior complexidade seguem para Fuerte Olimpo.
Cobertura privada continua recomendada para especialidades e urgências
de maior complexidade.
\newpage
\secaoDiagnostico{Sargento Jose Felix Lopez}{\mapaDistrito{0109}}{Sargento José Félix López (Puentesino) oferece um dos maiores potenciais
de isolamento do Norte paraguaio, mas exige um alto compromisso com a
segurança tática e autonomia logística. É um local de ``fronteira viva''
com desafios proporcionais ao seu isolamento.

\subsubsection{Por que sim}

\subsubsection{Por que não}

\subsubsection{Recomendação
Técnica}

Ideal para estabelecimentos que funcionam como postos avançados ou
refúgios táticos. Não recomendado para residência civil sem suporte de
segurança e logística independente. Referência atual de score: GSS 5.6.

\begin{itemize}
\tightlist
\item
  \begin{enumerate}
  \def\labelenumi{\arabic{enumi})}
  \tightlist
  \item
    Dados censitários (INE\footnote{\fineINE} 2022).
  \end{enumerate}
\item
  \begin{enumerate}
  \def\labelenumi{\arabic{enumi})}
  \setcounter{enumi}{1}
  \tightlist
  \item
    Relatórios de segurança regional (CODI 2025).
  \end{enumerate}
\item
  \begin{enumerate}
  \def\labelenumi{\arabic{enumi})}
  \setcounter{enumi}{2}
  \tightlist
  \item
    Informações sobre o Parque Nacional Paso Bravo (MADES).
  \end{enumerate}
\item
  \begin{enumerate}
  \def\labelenumi{\arabic{enumi})}
  \setcounter{enumi}{3}
  \tightlist
  \item
    Lei de Segurança Fronteiriça (2532/\allowbreak 05).
  \end{enumerate}
\end{itemize}}{dist:01_Concepción:Sargento_Jose_Felix_Lopez}
\begin{table}[ht!]
\small \begin{tabular}{p{3.0cm}p{0.8cm}p{6.0cm}} \toprule
\textbf{Critério} & \textbf{Nota} & \textbf{Justificativa} \\ \midrule
A - Ameaças Estratégicas & 4.0 & Zona de fronteira seca ativa e hotspot de grupos irregulares \\
B - Risco Social & 5.5 & Densidade baixíssima compensada pelo alto risco social/\allowbreak segurança regional \\
C - Riscos Naturais & 6.5 & Risco de isolamento hídrico moderado; geologia estável \\
D - Autossuficiência & 7.5 & Solo fértil e abundância hídrica natural; potencial de autossuficiência elevado \\
E - Institucional & 5.0 & Município recente; restrições legais rigorosas para estrangeiros \\
\midrule \textbf{GSS FINAL} & \textbf{5.6} & \textbf{Moderadamente Seguro (Recomendado apenas para perfis com experiência tática e foco em autossuficiência isolada).} \\ \bottomrule \end{tabular}
\end{table}
\subsubsection{Vulnerabilidades e
Mitigação}

\begin{itemize}
\tightlist
\item
  \textbf{Insegurança Regional:} Atividade insurgente e criminosa na
  fronteira seca (Mitigação: sistemas de segurança redundantes,
  comunicação via satélite e apoio à FTC).
\item
  \textbf{Isolamento Logístico:} Estradas de terra intransitáveis
  (Mitigação: estoque de suprimentos para 60 dias e posse de veículos de
  tração 4x4 pesada).
\item
  \textbf{Instabilidade Energética:} Rede elétrica pouco confiável
  (Mitigação: instalação de sistemas solares off-grid com baterias de
  lítio).
\item
  \textbf{Restrições de Propriedade:} Zona de fronteira exige
  estratégias jurídicas para investidores estrangeiros.
\end{itemize}

\subsubsection{Dossiê de Campo}
\begin{description}[style=nextline,leftmargin=0.5cm,font=\bfseries]
\item[Coordenadas]  22°31'S 56°52'W.
\item[Topografia]  Relevo ondulado com vastas áreas de pastagens e florestas. Proximidade imediata com o Parque Nacional Paso Bravo (grande reserva de biodiversidade e isolamento). Geologicamente estável.
\item[Alvos Estratégicos]  Grandes estâncias de pecuária de exportação, postos de fronteira seca com o Brasil (Bella Vista).
\item[Fallout]  Ventos predominantes do Norte/\allowbreak Nordeste.
\item[População]  6.153 (Censo 2022 Final).
\item[Densidade]  \textasciitilde 2,74 hab/\allowbreak km² (Baixa densidade, ideal para evitar grandes massas populacionais e garantir privacidade estratégica).
\item[Vias de Acesso]  Acesso Terrestre Crítico. Estradas de terra que sofrem severamente em períodos de chuva. Conexão principal com Concepción e Bella Vista Norte. Isolamento físico frequente.
\item[Serviços]  Infraestrutura básica de saúde (Puestos de Salud) e educação. Dependência crítica de Concepción para alta complexidade.
\item[Custo de Vida]  Reduzido para produtos locais, mas elevado para insumos industrializados devido ao frete.
\item[Preço da Terra]  US\$ 2.000-5.000/\allowbreak ha (áreas de fronteira pecuária e silvicultura).
\item[Hidrologia]  Risco de Isolamento por Chuvas. Transbordamento de arroios tributários do Rio Apa e estradas intransitáveis em períodos de chuva intensa são os principais riscos operacionais.
\item[Clima]  Tropical Úmido. Verões quentes e úmidos com tempestades severas.
\item[Energia]  Rede nacional (ANDE\footnote{\fineANDE}) instável por ser área remota e final de linha; interrupções frequentes.
\item[Água]  Poços artesianos e arroios locais perenes.
\item[Qualidade do Solo]  Solo com boa aptidão para pastagem e agricultura de subsistência; proximidade com áreas de reserva garante biodiversidade.
\item[Recursos Locais]  Pecuária de corte extensiva, exploração florestal sustentável e agricultura familiar. Alto potencial para autossuficiência básica.
\item[Segurança]  Zona de Alta Sensibilidade. Histórico de atuação de grupos insurgentes (EPP\footnote{\fineEPP}/\allowbreak ACA) e rotas de narcotráfico em fronteira seca. Exige alto nível de preparação tática e vigilância. Presença da FTC (Força de Tarefa Conjunta) é uma constante necessária.
\item[Leis Local: Município relativamente novo (2011) com forte presença de cultura de fronteira. Restrição de Fronteira (Lei 2532/\allowbreak 05)]  Proibição de compra de terras rurais por estrangeiros de países limítrofes na faixa de 50km.
\end{description}
\newpage
\secaoDiagnostico{Yby Yau}{\mapaDistrito{0107}}{Yby Yaú é o coração logístico do Norte paraguaio. Embora ofereça
recursos naturais e conectividade excepcionais, sua pontuação GSS de 5.4
é pesadamente penalizada pela insegurança crônica e pelo risco de ser um
alvo tático em qualquer cenário de instabilidade nacional.

\subsubsection{Por que sim}

\subsubsection{Por que não}

\subsubsection{Recomendação
Técnica}

Ideal para operações agroindustriais de larga escala que possam arcar
com custos elevados de segurança privada. Não recomendado como refúgio
familiar passivo devido à ``visibilidade'' do local. Referencia atual de
score: GSS 5.4.

\begin{itemize}
\tightlist
\item
  \begin{enumerate}
  \def\labelenumi{\arabic{enumi})}
  \tightlist
  \item
    Dados censitários (INE\footnote{\fineINE} 2022).
  \end{enumerate}
\item
  \begin{enumerate}
  \def\labelenumi{\arabic{enumi})}
  \setcounter{enumi}{1}
  \tightlist
  \item
    Relatórios de segurança regional (CODI 2025).
  \end{enumerate}
\item
  \begin{enumerate}
  \def\labelenumi{\arabic{enumi})}
  \setcounter{enumi}{2}
  \tightlist
  \item
    Mapa de rodovias e logística (MOPC 2026).
  \end{enumerate}
\item
  \begin{enumerate}
  \def\labelenumi{\arabic{enumi})}
  \setcounter{enumi}{3}
  \tightlist
  \item
    Jurisprudência sobre a Lei de Fronteira (2532/05).
  \end{enumerate}
\end{itemize}}{dist:01_Concepcion:Yby_Yau}
\begin{table}[ht!]
\small \begin{tabular}{p{3.0cm}p{0.8cm}p{6.0cm}} \toprule
\textbf{Critério} & \textbf{Nota} & \textbf{Justificativa} \\ \midrule
A - Ameaças Estratégicas & 3.5 & Nó logístico vital; alvo de controle militar e civil em crises \\
B - Risco Social & 4.0 & Alta densidade de criminalidade organizada e insurgência regional \\
C - Riscos Naturais & 7.0 & Baixo risco de desastres naturais; geologia muito estável \\
D - Autossuficiência & 7.5 & Solo fértil, abundância de recursos hídricos e proteína animal \\
E - Institucional & 6.0 & Fora da restrição de 50km da fronteira; ambiente institucional desafiado pela violência \\
\midrule \textbf{GSS FINAL} & \textbf{5.4} & \textbf{Moderadamente Seguro (Recomendado apenas para perfis táticos com interesse comercial/agroindustrial).} \\ \bottomrule \end{tabular}
\end{table}
\subsubsection{Vulnerabilidades e
Mitigação}

\begin{itemize}
\tightlist
\item
  \textbf{Insegurança Pública:} Alta incidência de crimes violentos
  ligados ao tráfico e insurgência (Mitigação: residência fora do núcleo
  urbano imediato, sistemas de monitoramento e blindagem patrimonial).
\item
  \textbf{Alvo Logístico:} Risco de bloqueios rodoviários e requisição
  de suprimentos em crises (Mitigação: manter rotas de fuga rurais não
  pavimentadas conhecidas).
\item
  \textbf{Dependência de Saúde:} Distância de centros cirúrgicos
  (Mitigação: treinamento avançado de primeiros socorros e kit trauma
  completo).
\end{itemize}

\subsubsection{Dossiê de Campo}
\begin{description}[style=nextline,leftmargin=0.5cm,font=\bfseries]
\item[Coordenadas]  23°26'S 56°29'W.
\item[Topografia]  Relevo ondulado, zona de transição para a Cordilheira de Amambay. Altitude \textasciitilde 200m. Geologicamente estável, risco sísmico desprezível.
\item[Alvos Estratégicos]  Entroncamento das Rutas PY05 (Leste-Oeste) e PY08 (Norte-Sul) e proximidade com a PY03. Ponto vital absoluto para a logística do Norte do Paraguai. Hub comercial e de trânsito de gado e grãos.
\item[Fallout]  Ventos predominantes do Norte/Nordeste (quentes) e Sul (inverno).
\item[População]  19.468 (Censo 2022 Final).
\item[Densidade]  \textasciitilde 8,04 hab/km² (Moderada, concentrada no núcleo urbano do cruzamento rodoviário).
\item[Vias de Acesso]  Nó Rodoviário Crítico. Excelente conectividade em tempos de paz, mas risco extremo de controle por grupos armados ou forças militares em caso de conflito. Gargalo para qualquer movimentação Norte-Sul.
\item[Serviços]  Comércio altamente desenvolvido, bancos e centros de saúde básicos. Dependência de Concepción ou Pedro Juan Caballero (ambos a \textasciitilde 100 km) para alta complexidade.
\item[Custo de Vida]  Baixo; produção local de carne e hortifrúti favorece a cesta básica.
\item[Preço da Terra]  US\$ 1.200-2.000/ha (terras agrícolas e pastagens).
\item[Hidrologia]  Diversos arroios cruzam o distrito. Baixo risco de inundações em comparação com distritos ribeirinhos. Saneamento básico em fase de expansão.
\item[Clima]  Subtropical Úmido. Precipitação anual de 1.400-1.600 mm.
\item[Energia]  Rede nacional (Itaipu). Importante subestação de distribuição regional.
\item[Água]  Poços artesianos e arroios locais perenes.
\item[Qualidade do Solo]  Latossolos vermelhos e solos arenosos; profundos e mecanizáveis, exigem correção para alta produtividade.
\item[Recursos Locais]  Polo pecuário intensivo e produção de grãos (soja/milho). Elevadíssimo potencial para autossuficiência proteica e energética (se houver autonomia solar).
\item[Leis Local]  Município fora da restrição absoluta da Lei de Fronteira (distância > 50km da linha seca), facilitando a titularidade para estrangeiros.
\end{description}
\textbf{Fonte climática:} NASA POWER Climatology API (período 2001-2020)

\textbf{Fonte luminosa:} estimativa world atlas

\paragraph{Irradiação Solar
(kWh/m²/dia)}

{\def\LTcaptype{none} % do not increment counter
\begin{longtable}[]{@{}
  >{\raggedright\arraybackslash}p{(\linewidth - 24\tabcolsep) * \real{0.0746}}
  >{\raggedright\arraybackslash}p{(\linewidth - 24\tabcolsep) * \real{0.0746}}
  >{\raggedright\arraybackslash}p{(\linewidth - 24\tabcolsep) * \real{0.0746}}
  >{\raggedright\arraybackslash}p{(\linewidth - 24\tabcolsep) * \real{0.0746}}
  >{\raggedright\arraybackslash}p{(\linewidth - 24\tabcolsep) * \real{0.0746}}
  >{\raggedright\arraybackslash}p{(\linewidth - 24\tabcolsep) * \real{0.0746}}
  >{\raggedright\arraybackslash}p{(\linewidth - 24\tabcolsep) * \real{0.0746}}
  >{\raggedright\arraybackslash}p{(\linewidth - 24\tabcolsep) * \real{0.0746}}
  >{\raggedright\arraybackslash}p{(\linewidth - 24\tabcolsep) * \real{0.0746}}
  >{\raggedright\arraybackslash}p{(\linewidth - 24\tabcolsep) * \real{0.0746}}
  >{\raggedright\arraybackslash}p{(\linewidth - 24\tabcolsep) * \real{0.0746}}
  >{\raggedright\arraybackslash}p{(\linewidth - 24\tabcolsep) * \real{0.0746}}
  >{\raggedright\arraybackslash}p{(\linewidth - 24\tabcolsep) * \real{0.1045}}@{}}
\toprule\noalign{}
\begin{minipage}[b]{\linewidth}\raggedright
Jan
\end{minipage} & \begin{minipage}[b]{\linewidth}\raggedright
Fev
\end{minipage} & \begin{minipage}[b]{\linewidth}\raggedright
Mar
\end{minipage} & \begin{minipage}[b]{\linewidth}\raggedright
Abr
\end{minipage} & \begin{minipage}[b]{\linewidth}\raggedright
Mai
\end{minipage} & \begin{minipage}[b]{\linewidth}\raggedright
Jun
\end{minipage} & \begin{minipage}[b]{\linewidth}\raggedright
Jul
\end{minipage} & \begin{minipage}[b]{\linewidth}\raggedright
Ago
\end{minipage} & \begin{minipage}[b]{\linewidth}\raggedright
Set
\end{minipage} & \begin{minipage}[b]{\linewidth}\raggedright
Out
\end{minipage} & \begin{minipage}[b]{\linewidth}\raggedright
Nov
\end{minipage} & \begin{minipage}[b]{\linewidth}\raggedright
Dez
\end{minipage} & \begin{minipage}[b]{\linewidth}\raggedright
Média
\end{minipage} \\
\midrule\noalign{}
\endhead
\bottomrule\noalign{}
\endlastfoot
6.44 & 5.92 & 5.66 & 4.74 & 3.62 & 3.27 & 3.51 & 4.32 & 4.78 & 5.46 &
6.22 & 6.37 & \textbf{5.03} \\
\end{longtable}
}

\textbf{Inclinação solar recomendada:} 23° N (anual) · 33° N (inverno
jun-ago) · 13° N (verão nov-jan)

\paragraph{Precipitação (mm/mês)}

{\def\LTcaptype{none} % do not increment counter
\begin{longtable}[]{@{}
  >{\raggedright\arraybackslash}p{(\linewidth - 24\tabcolsep) * \real{0.0704}}
  >{\raggedright\arraybackslash}p{(\linewidth - 24\tabcolsep) * \real{0.0704}}
  >{\raggedright\arraybackslash}p{(\linewidth - 24\tabcolsep) * \real{0.0704}}
  >{\raggedright\arraybackslash}p{(\linewidth - 24\tabcolsep) * \real{0.0704}}
  >{\raggedright\arraybackslash}p{(\linewidth - 24\tabcolsep) * \real{0.0704}}
  >{\raggedright\arraybackslash}p{(\linewidth - 24\tabcolsep) * \real{0.0704}}
  >{\raggedright\arraybackslash}p{(\linewidth - 24\tabcolsep) * \real{0.0704}}
  >{\raggedright\arraybackslash}p{(\linewidth - 24\tabcolsep) * \real{0.0704}}
  >{\raggedright\arraybackslash}p{(\linewidth - 24\tabcolsep) * \real{0.0704}}
  >{\raggedright\arraybackslash}p{(\linewidth - 24\tabcolsep) * \real{0.0704}}
  >{\raggedright\arraybackslash}p{(\linewidth - 24\tabcolsep) * \real{0.0704}}
  >{\raggedright\arraybackslash}p{(\linewidth - 24\tabcolsep) * \real{0.0704}}
  >{\raggedright\arraybackslash}p{(\linewidth - 24\tabcolsep) * \real{0.1549}}@{}}
\toprule\noalign{}
\begin{minipage}[b]{\linewidth}\raggedright
Jan
\end{minipage} & \begin{minipage}[b]{\linewidth}\raggedright
Fev
\end{minipage} & \begin{minipage}[b]{\linewidth}\raggedright
Mar
\end{minipage} & \begin{minipage}[b]{\linewidth}\raggedright
Abr
\end{minipage} & \begin{minipage}[b]{\linewidth}\raggedright
Mai
\end{minipage} & \begin{minipage}[b]{\linewidth}\raggedright
Jun
\end{minipage} & \begin{minipage}[b]{\linewidth}\raggedright
Jul
\end{minipage} & \begin{minipage}[b]{\linewidth}\raggedright
Ago
\end{minipage} & \begin{minipage}[b]{\linewidth}\raggedright
Set
\end{minipage} & \begin{minipage}[b]{\linewidth}\raggedright
Out
\end{minipage} & \begin{minipage}[b]{\linewidth}\raggedright
Nov
\end{minipage} & \begin{minipage}[b]{\linewidth}\raggedright
Dez
\end{minipage} & \begin{minipage}[b]{\linewidth}\raggedright
Total/ano
\end{minipage} \\
\midrule\noalign{}
\endhead
\bottomrule\noalign{}
\endlastfoot
178 & 183 & 118 & 133 & 121 & 80 & 50 & 43 & 83 & 159 & 190 & 174 &
\textbf{1517 mm} \\
\end{longtable}
}

\paragraph{Poluição Luminosa}

{\def\LTcaptype{none} % do not increment counter
\begin{longtable}[]{@{}ll@{}}
\toprule\noalign{}
Parâmetro & Valor \\
\midrule\noalign{}
\endhead
\bottomrule\noalign{}
\endlastfoot
Escala Bortle & 3 --- Céu rural \\
Radiância artificial & 1.5 nW/cm²/sr \\
\end{longtable}
}
\paragraph{Solo (SoilGrids 2.0, média ponderada 0--30
cm)}

{\def\LTcaptype{none} % do not increment counter
\begin{longtable}[]{@{}ll@{}}
\toprule\noalign{}
Parâmetro & Valor \\
\midrule\noalign{}
\endhead
\bottomrule\noalign{}
\endlastfoot
pH (H$_2$O) & 5.5 \\
Carbono orgânico (SOC) & 18.06 g/kg \\
Argila & 21.73 \% \\
Areia & 60.38 \% \\
Silte (calc.) & 17.9 \% \\
Densidade aparente & 1.25 g/cm³ \\
\textbf{Aptidão agrícola} & \textbf{Alta} \\
\end{longtable}
}

Fonte: ISRIC SoilGrids 2.0 via WCS. Coords: -23.4333°, -56.4833°. Média
ponderada camadas 0-5, 5-15, 15-30 cm.
\subsubsection{Indicadores Sociais}

\textbf{Fontes:} PNUD 2020, INE\footnote{\fineINE} Censo 2022, INE\footnote{\fineINE} EPHC 2023, Ministerio
Público 2024\\
\textbf{Nota:} Valores marcados como \emph{(dept.)} referem-se ao
departamento de Concepción; valores \emph{(dist.)} são específicos deste
distrito. Dados marcados \emph{(est.)} são estimativas.

{\def\LTcaptype{none} % do not increment counter
\begin{longtable}[]{@{}
  >{\raggedright\arraybackslash}p{(\linewidth - 4\tabcolsep) * \real{0.4231}}
  >{\raggedright\arraybackslash}p{(\linewidth - 4\tabcolsep) * \real{0.2692}}
  >{\raggedright\arraybackslash}p{(\linewidth - 4\tabcolsep) * \real{0.3077}}@{}}
\toprule\noalign{}
\begin{minipage}[b]{\linewidth}\raggedright
Indicador
\end{minipage} & \begin{minipage}[b]{\linewidth}\raggedright
Valor
\end{minipage} & \begin{minipage}[b]{\linewidth}\raggedright
Âmbito
\end{minipage} \\
\midrule\noalign{}
\endhead
\bottomrule\noalign{}
\endlastfoot
IDH (2020) & 0.650 & dist. (est.) \\
Ranking IDH nacional & 9/18 & dept. \\
Esperança de vida & 72,4 anos & dept. \\
Escolaridade média & 8.0 anos & dist. \\
RNB per capita (USD PPA) & 8.400 & dept. \\
Pobreza monetária (\%) & 31,1\% & dept. \\
Pobreza extrema (\%) & 7,8\% & dept. \\
Índice de Gini & 0,460 & dept. \\
Acesso a água potável (\%) & 74,2\% & dept. \\
Acesso a saneamento (\%) & 71,8\% & dept. \\
Taxa de homicídios (est., /100k hab) & \textasciitilde12--15
(estimativa, acima da média nacional) & dept. \\
Índice de segurança & baixo & dept. \\
População (Censo 2022) & 19.468 hab. & dist. \\
Idade mediana & 26.0 anos & dist. \\
Taxa de fecundidade & 2 & dist. \\
\end{longtable}
}
\subsubsection{Combustível}

\textbf{Referência:} PETROPAR / postos locais (2024) \textbf{Tipo de
localidade:} Interior

{\def\LTcaptype{none} % do not increment counter
\begin{longtable}[]{@{}lll@{}}
\toprule\noalign{}
Combustível & USD/litro & Gs/litro (aprox.) \\
\midrule\noalign{}
\endhead
\bottomrule\noalign{}
\endlastfoot
Gasolina 93 oct & 1.01 & 7,474 \\
Gasolina 97 oct (premium) & 1.11 & 8,214 \\
Diesel & 0.94 & 6,956 \\
\end{longtable}
}

\begin{quote}
Preços podem variar ±5\% conforme posto e sazonalidade. Chaco e interior
remoto apresentam maior variação.
\end{quote}

\subsubsection{Cobertura Celular}

\textbf{Fonte:} CONATEL PY / operadoras (2024)

{\def\LTcaptype{none} % do not increment counter
\begin{longtable}[]{@{}ll@{}}
\toprule\noalign{}
Parâmetro & Valor \\
\midrule\noalign{}
\endhead
\bottomrule\noalign{}
\endlastfoot
Cobertura 4G população (dept.) & 78\% \\
Cobertura 4G área rural & 55\% \\
Melhor operadora & Tigo \\
Qualidade rural & moderada \\
\end{longtable}
}

\begin{quote}
Para áreas rurais fora do núcleo urbano, recomenda-se chip Tigo como
principal e Personal como backup.
\end{quote}

\subsubsection{Internet}

\textbf{Fonte:} CONATEL / Speedtest Ookla (2024)

{\def\LTcaptype{none} % do not increment counter
\begin{longtable}[]{@{}ll@{}}
\toprule\noalign{}
Parâmetro & Valor \\
\midrule\noalign{}
\endhead
\bottomrule\noalign{}
\endlastfoot
Velocidade média download & 38 Mbps \\
Domicílios com internet (dept.) & 48\% \\
Tecnologia predominante & rádio \\
Opção rural & Starlink disponível (\textasciitilde USD 44/mês) \\
\end{longtable}
}

\subsubsection{Mercado Imobiliário e Terra
Rural}

\textbf{Fonte:} INDERT / Clasificados.com.py (2024)

{\def\LTcaptype{none} % do not increment counter
\begin{longtable}[]{@{}ll@{}}
\toprule\noalign{}
Tipo & Referência \\
\midrule\noalign{}
\endhead
\bottomrule\noalign{}
\endlastfoot
Terra agrícola alta prod. (USD/ha) & 3,500 \\
Imóvel urbano (USD/m²) & 650 \\
Aluguel 2 quartos (USD/mês) & 260 \\
\end{longtable}
}

\begin{quote}
Valores de referência departamental. Localidades menores podem ter
preços 20--40\% abaixo da capital departamental.
\end{quote}

\subsubsection{Saúde}

\textbf{Fonte:} MSPBS / IPS Paraguay (2024-2026), consolidação
departamental e proxy local

{\def\LTcaptype{none} % do not increment counter
\begin{longtable}[]{@{}ll@{}}
\toprule\noalign{}
Serviço & Disponibilidade \\
\midrule\noalign{}
\endhead
\bottomrule\noalign{}
\endlastfoot
USF / Posto de Saúde & sim \\
Hospital Regional & não \\
IPS (seguro social) & sim \\
Farmácia & sim \\
Distância ao hospital de referência & \textasciitilde130 km
(Concepcion) \\
\end{longtable}
}

\textbf{Principais estabelecimentos:} USF local; referencia hospitalar
em Concepcion

\textbf{Observação para imigrantes:} Atendimento primario local ou em
raio curto; casos de maior complexidade seguem para Concepcion.
Cobertura privada continua recomendada para especialidades e urgências
de maior complexidade.
